% !TEX root = ../print.tex
% Chapter 4 == section 4 of the article (notes/PLAN-article-structure.md, "The outline"): the
% (G, W) axiom scheme and the classical affine theory, G = GL(d) and W the Loewner order. Sections:
% 4.1 the scheme and the passage from operators to kernels (sec:affine-covariance), 4.2 the parameter
% space (sec:parameter-space), 4.3 the null array (sec:null-array), 4.4 the single step and rigidity
% (sec:affine-single-step). The article's 4.5 (prior art: Kelker's lemma, Lindeberg's generalized
% axiomatics) is prose and has no node. Until Q-0278 (2026-10-04) these nodes were the first section
% of 02-affine.tex, which held the draft's sections 2.2-2.4; no label, statement or proof changed in
% the move (scripts/check-statements-unchanged.py). Provenance of the draft's nodes: Q-0152 (hub
% queue).
%
% Design, and where it departs from the draft. Every node here but one, and every family node of
% chapter 5, is stated at the level of KERNELS, not of operators. The one is lem:operators-to-kernels (Q-0207,
% on the author's answer (c) to Q-0205): it takes the draft's operator axioms, of Definition 2.4 and
% of Definition 2.1 alike, to the kernel definitions, by [V, Lem. 4.1]'s proof redone in R^d with
% balls for intervals, as notes/PLAN-d-dimensional-port.md S1 found possible; no later node reads
% it. None of the nodes reads the draft's Theorem 2.2 (the isotropic classification in R^d). That
% theorem now has a node, thm:isotropic-classification in 03-isotropic.tex (Q-0208), proved by
% reduction to [V, Thm. 7.3] through the one-dimensional marginal, with its polar reading
% prop:isotropic-polar-reading; the affine nodes were written before either existed and do not use
% them, and nothing here needs to. Their annotations named neither node and called the
% classification an unported claim; Q-0211 rewrote those annotations to name
% thm:isotropic-classification (and prop:isotropic-polar-reading where the explicit polar form (2.1)
% is meant) and to say, for each affine node, that its proof does not use them. The draft's Lemma
% 2.6 obtained the one-function form (2.2) through Theorem 2.2 and the negative definiteness of the
% increments through the null-array argument of [V, Thm. 5.2] ported to R^d; neither is needed for
% the rigidity. The one-function form follows from covariance alone (lem:affine-reduction), and the
% rigidity follows from ONE rank-deficient step whose increment is merely a probability measure
% (thm:rank-deficient-step, the draft's Prop. 2.14), through the probability-measure form of the
% draft's Lemma 2.7 (lem:charfn-one-on-subspace). So the draft's Theorem 2.8 is proved here as a
% corollary of its Proposition 2.14, and the order of the two is reversed. Nothing in the draft's
% sections 2.2-2.4 cites a source.
%
% Q-0265 (notes/PLAN-article-structure.md, items D1, D2, D4) added: def:gw-cascade-family, the
% (G, W) scheme on operators and kernels, of which def:affine-cascade-family is the kernel form at
% (GL(d), Loewner), with lem:operators-to-kernels(4) generalized to it; lem:parameter-space, why the
% parameter space is PD(d) and that the isotropic member is an axiom; thm:cascade-infinitely-divisible,
% the null array (ledger A6, Khintchine's theorem, admitted by the author 2026-10-03): along a W-chain
% every increment is infinitely divisible; and prop:similarity-cascade-families (now in
% 05-smaller-groups.tex), which discharges the "infinitely divisible increments" hypothesis of the A'
% nodes for (similarities, Loewner)-cascade families. The rigidity nodes still do not use infinite
% divisibility. Its clause (3), the orders with the path condition, is split to
% lem:path-condition-orders (Q-0281), which \uses lem:stable-step-order of chapter 5 (that the stable
% step cone's relation is a partial order, the order of prop:stable-step-cone(2), split off there);
% thm:cascade-infinitely-divisible \uses both. The forward edges are the article's order, not a cycle.
% CHANGED Q-0281: blind review R5-3 (this comment named prop:stable-step-cone and clause (3) before
% the two splits).
%
% Q-0280 added prop:isotropic-corner (section 4.1): the slice {sI} of the scheme's isotropic corner
% is def:isotropic-kernel-family, up to the variance gauge, which [V, Prop. 6.6] supplies; chapter
% 3's definition is kept as it stands and is the corner's slice. It reads three nodes of section
% 5.6, forward edges like thm:cascade-infinitely-divisible's.
%
% Q-0287 (the article's section 4 transcribed) changed two statements, in the safe direction and
% without changing what either claims: lem:operators-to-kernels(4) and (5) no longer name "the
% draft's Definition 2.4/2.1" (in the paper those numbers point at other statements), the axioms
% being listed in the clause as before; and thm:cascade-infinitely-divisible(3)(c) states that its
% three relations are partial orders without the justification "the first two directly and the
% third by prop:stable-step-cone(2)", which moved to the proof of (3) (the paper's section 4 comes
% before the stable step cone of its section 5). The annotation of thm:affine-rigidity carries the
% page-checked correction of the Skitovich anchor (Thm. 2, p. 190).
%
% Q-0289 (the blind review of the article's section 4, records/affine/REVIEW-paper-section-4.md)
% changed two statements, in the safe direction and without changing what either claims:
% def:gw-cascade-family and lem:operators-to-kernels(1) call the single-operator axioms of [V, Def.
% 3.1] they list "the single-operator axioms other than the reflection axiom (A3)" instead of "the
% four axioms that concern a single operator" (there are five: (A3) concerns one operator too); and
% thm:cascade-infinitely-divisible(3)(c) is split into three sentences (the hypotheses on K, the
% partial-order hypothesis, the chain), with the same quantifiers. thm:affine-rigidity's annotation
% now says which half on operators is not claimed.
%
% Stays prose, and why: the draft's Corollary 2.9 (its content is thm:affine-rigidity; the rest
% is a comparison with Lindeberg 2011, now read at the page and accurate --
% records/affine/READING-lindeberg2011generalized.md -- but still a comparison no node's proof
% depends on, so no ledger entry); it is the article's 4.5.

\chapter{The axiom scheme and the classical affine theory}
\label{ch:classical}

This chapter is the article's \S 4, and in the draft's numbering its \S 2.2. Its theorem is that in
dimension $d \ge 2$ covariance under the whole affine group, together with a cascade along the
Loewner order of covariance matrices, leaves the Gaussian alone (\S\ref{sec:affine-single-step}).
It opens with what is not in the draft: the axiom scheme of which the draft's affine and isotropic
definitions are corners, with the passage from operators to kernels (\S\ref{sec:affine-covariance});
why the parameter space is $\PD(d)$ (\S\ref{sec:parameter-space}); and the infinite divisibility of
the increments of every cascade along an order with the path condition, by Khintchine's null-array
theorem (\S\ref{sec:null-array}). The rigidity that follows uses none of the last two.

\section{The $(G, W)$ scheme, from operators to kernels}
\label{sec:affine-covariance}

The axiom scheme is stated once on operators and once on kernels (\ref{def:gw-cascade-family}); its
corner at the general linear group and the Loewner order is the draft's affine definition, on kernels
(\ref{def:affine-cascade-family}). The nodes are stated on kernels, one lemma
(\ref{lem:operators-to-kernels}) connecting them to the draft's operators. Read on the multiples of
$I$, the corner at the similarities and the rays is the isotropic theory of chapter~3, up to a change
of the scale coordinate (\ref{prop:isotropic-corner}).

% not in the draft: notes/PLAN-article-structure.md, development item D2 (the author's decisions
% of 2026-10-03; queue Q-0265). The one axiom scheme the three main sections are corners of, stated
% once on operators and once on kernels. The draft's Definitions 2.1 and 2.4 are two of its corners.
% What writing it down found: indexing by PSD(d) with the action A.Sigma = A Sigma A^T already
% contains the axiom that the family has an isotropic member (lem:parameter-space).
% CHANGED Q-0281: typed (the statement skeleton). The definition's displayed claims -- (0, I) in P_W,
% the Loewner and ray orders as invariant partial orders containing the rays, and the transform of an
% image measure -- are split to lem:gw-scheme-claims (ADR-0025; records/formalization/PRICING.md,
% finding 4), a Lean definition being unable to carry a claim; the text below is unchanged, and the
% \lean tag names the claim lemma's declarations too, so that it covers every sentence.
% Kept \notready although every tagged declaration is now defined or proved (wave 1A, blind review
% W1A-2): the operator form encodes D_A relationally, its equivalence with the operator identity
% being a docstring of Defs/Operators.lean only (the author's open call 1,
% records/formalization/SKELETON.md). The list of corners also asserts that (S(d), rays) and
% (S(d), Loewner) are wedge data, which needs S(d) closed in GL(d); that is now proved,
% AffineHemigroup.simSet_isClosedSubgroupGL (wave 2, group C, in the proof of
% prop:isotropic-corner(4)), but not tagged here while the operator-form reason stands.
% CHANGED Q-0281 wave 2: blind review W2C-8 (the closedness of S(d) is proved; the operator-form
% reason for \notready still holds).
% CHANGED Q-0281 wave M, phase 1: both reasons for \notready above are now settled. The closedness
% of S(d) is proved (simSet_isClosedSubgroupGL, as said). The author's open call 1 is closed (the
% author's decision of 2026-10-07): the operator form is re-typed literally, as equations of
% operators on L^1(R^d) with the shared library's T_a = ScaleSpace.transL1E a and
% D_A = ScaleSpace.dilL1E (ScaleSpace.matEquiv A hA), normalised as here, |det A|^{-1} f(A^{-1} .)
% (AffineHemigroup.Defs.Operators); the covariance is the conjugation D_A Phi D_A^{-1} as written.
% The definition AffineHemigroup.Intertwines, the relational reading, is gone from the tag: that
% reading is now a theorem, AffineHemigroup.covariant_iff_relational and its siblings in
% AffineHemigroup.OperatorForm. The node stays \notready only until phase 2 proves its claims
% against the new form (and tags simSet_isClosedSubgroupGL for the corners); it goes green then.
% The text is unchanged; only the Lean changed.
% CHANGED Q-0281 wave M, phase 2: green, as the comment above says it goes. The tag gains
% AffineHemigroup.simSet_isClosedSubgroupGL (the closedness of S(d) that the corners with the
% similarities need); every tagged declaration is defined or proved, the shared library's
% ScaleSpace.* among them, and the operator form's claims are proved against the literal form
% (lem:operators-to-kernels, lem:kernels-to-operators). The text is unchanged.
% CHANGED Q-0469 (wording only; records/affine/GATE-affine.md, Q-0317 row "§3 18"): the ray order is
% described in words, "under which Sigma precedes Sigma'", where it was written with the generic
% \le_W of an arbitrary datum; "read on the slice {sI : s >= 0}" is "read on the multiples sI,
% s >= 0, of the identity", "slice" being defined only in prop:isotropic-corner. No claim changes.
% CHANGED pass 6 of the affine paper (layout only; the author's review of section 3.1, 2026-10-27,
% records/affine/REVIEW-author-paper-affine.md): the operator axioms are one per line, as the kernel
% axioms are, with (A1), (A2), (A4), (A5) written out and the sentence citing the line paper's
% single-operator axioms after the list; (ND) gains its name; covariance is labelled (Cov); the four corners are a list, and the
% classical one is named "the classical case". No axiom, claim or Lean tag changes.
\begin{definition}[the $(G, W)$ cascade scheme]
\label{def:gw-cascade-family}
\lean{AffineHemigroup.IsWedgeDatum, AffineHemigroup.wedgePairs,
AffineHemigroup.IsGWOperatorFamily, AffineHemigroup.IsSingleOperator,
AffineHemigroup.IsGWCascadeFamily, AffineHemigroup.loewnerLE,
AffineHemigroup.rayLE, AffineHemigroup.push, AffineHemigroup.act, AffineHemigroup.glSet,
AffineHemigroup.simSet, AffineHemigroup.wedge_zero_le, AffineHemigroup.wedge_zero_one_mem,
AffineHemigroup.isWedgeDatum_loewner, AffineHemigroup.isWedgeDatum_ray,
AffineHemigroup.glSet_isClosedSubgroupGL, AffineHemigroup.charFun_push,
AffineHemigroup.simSet_isClosedSubgroupGL,
ScaleSpace.transL1E, ScaleSpace.dilL1E, ScaleSpace.matEquiv,
ScaleSpace.dilL1E_matEquiv_comp}\leanok
Fix an integer $d \ge 1$. Write $\PSD(d)$ for the real symmetric positive semidefinite
$d \times d$ matrices and $\PD(d) \subseteq \PSD(d)$ for the positive definite ones. The group
$\GL(d)$ acts on $\PSD(d)$ by $A \cdot \Sigma := A\Sigma A^{\tp}$, on probability measures on
$\RR^d$ by pushforward, $A_*\mu := \mu \circ A^{-1}$, and on $L^1(\RR^d)$ by
$D_Af := |\det A|^{-1}f(A^{-1}\,\cdot)$; with
$\widehat\mu(\omega) = \int e^{i\langle\omega,x\rangle}\mu(dx)$ one has
$\widehat{A_*\mu}(\omega) = \widehat\mu(A^{\tp}\omega)$. A \emph{wedge datum in dimension $d$} is a
pair $(G, W)$ of
\begin{itemize}
\item a closed subgroup $G$ of $\GL(d)$ containing the dilations $\lambda I$, $\lambda > 0$, and
\item a partial order $\le_W$ on $\PSD(d)$ that is \emph{$G$-invariant}, $\Sigma \le_W \Sigma'$
implying $A \cdot \Sigma \le_W A \cdot \Sigma'$ for every $A \in G$, and \emph{contains the rays},
$c\,\Sigma \le_W c'\,\Sigma$ whenever $\Sigma \in \PSD(d)$ and $0 \le c \le c'$.
\end{itemize}
Write $P_W$ for the set of pairs $(\Sigma, \Sigma')$ with $\Sigma \le_W \Sigma'$, with the topology
of $\PSD(d) \times \PSD(d)$. Since $W$ contains the rays, $0 \le_W \Sigma$ for every $\Sigma$, and
in particular $(0, I) \in P_W$.
\begin{enumerate}
\item \emph{(Operator form.)} A \emph{$(G, W)$-cascade family of operators} is a family
$(\Phi_{\Sigma,\Sigma'})_{(\Sigma,\Sigma') \in P_W}$ of linear operators on the real space
$L^1(\RR^d)$ such that
\begin{enumerate}
\item[(A1)] each $\Phi_{\Sigma,\Sigma'}$ is a bounded operator on $L^1(\RR^d)$;
\item[(A2)] each $\Phi_{\Sigma,\Sigma'}$ commutes with the translations $T_af := f(\cdot - a)$,
$a \in \RR^d$;
\item[(A4)] each $\Phi_{\Sigma,\Sigma'}$ maps nonnegative functions to nonnegative functions;
\item[(A5)] each $\Phi_{\Sigma,\Sigma'}$ preserves $\int f$ for $f \ge 0$;
\item[(A6)] \emph{(cascade)} $\Phi_{\Sigma,\Sigma} = \mathrm{Id}$ and
$\Phi_{\Sigma',\Sigma''}\Phi_{\Sigma,\Sigma'} = \Phi_{\Sigma,\Sigma''}$ whenever
$\Sigma \le_W \Sigma' \le_W \Sigma''$;
\item[(A7)] \emph{(continuity)} for every $f \in L^1(\RR^d)$ the map $(\Sigma,\Sigma') \mapsto \Phi_{\Sigma,\Sigma'}f$
is continuous from $P_W$ into $L^1(\RR^d)$;
\item[(ND)] \emph{(nondegeneracy)} $\Phi_{0,I} \ne \mathrm{Id}$;
\item[(Cov)] \emph{(covariance)} $D_A\Phi_{\Sigma,\Sigma'}D_A^{-1} = \Phi_{A\cdot\Sigma,\,A\cdot\Sigma'}$ for every
$A \in G$ and every $(\Sigma,\Sigma') \in P_W$.
\end{enumerate}
Axioms (A1), (A2), (A4) and (A5) are the single-operator axioms of
\restated{V}{Def.~3.1}{def:line-family} other than the reflection axiom (A3), read in $\RR^d$.
\item \emph{(Kernel form.)} A \emph{$(G, W)$-cascade family of kernels} is a family
$(\mu_{\Sigma,\Sigma'})_{(\Sigma,\Sigma') \in P_W}$ of probability measures on $\RR^d$ such that
\begin{enumerate}
\item[(K1)] \emph{(cascade)} $\mu_{\Sigma,\Sigma} = \delta_0$, and
$\mu_{\Sigma,\Sigma''} = \mu_{\Sigma',\Sigma''} * \mu_{\Sigma,\Sigma'}$ whenever
$\Sigma \le_W \Sigma' \le_W \Sigma''$;
\item[(K2)] \emph{(continuity)} for every $\omega \in \RR^d$ the map
$(\Sigma,\Sigma') \mapsto \widehat\mu_{\Sigma,\Sigma'}(\omega)$ is continuous on $P_W$;
\item[(K3)] \emph{($G$-covariance)} $A_*\mu_{\Sigma,\Sigma'} = \mu_{A\cdot\Sigma,\,A\cdot\Sigma'}$ for
every $A \in G$ and every $(\Sigma,\Sigma') \in P_W$;
\item[(K4)] \emph{(nondegeneracy)} $\mu_{0,I} \ne \delta_0$.
\end{enumerate}
We write $\mu_\Sigma := \mu_{0,\Sigma}$ for the \emph{kernel from scale zero}.
\end{enumerate}
The \emph{Loewner order}, $\Sigma \le \Sigma'$ when $\Sigma' - \Sigma \in \PSD(d)$, and the
\emph{ray order}, under which $\Sigma$ precedes $\Sigma'$ when $\Sigma = 0$ or $\Sigma' = c\,\Sigma$
for some $c \ge 1$, contain the rays and are invariant under every subgroup of $\GL(d)$. The
\emph{corners} of the scheme are these four data:
\begin{itemize}
\item the similarities $\{cQ : c > 0,\ Q \in \Orth(d)\}$ with the ray order, read on the multiples
$sI$, $s \ge 0$, of the identity (the isotropic theory);
\item $(\GL(d), \text{Loewner})$, the classical case, whose kernel form is
\ref{def:affine-cascade-family};
\item $(\GL(d), \text{rays})$;
\item the similarities with the Loewner order.
\end{itemize}
\statusT\quad\emph{(Vocabulary, and no claim. Provenance: the plan's item D2,
\texttt{notes/PLAN-article-structure.md}, \S\ ``The organizing idea''. The operator form is the
draft's Definition~2.4 with $\GL(d)$ and the Loewner order replaced by $(G, W)$, and the passage to
the kernel form is \ref{lem:operators-to-kernels}(4), whose proof is per operator and does not see
$(G, W)$. \textbf{What writing it down found.} Indexing the family by $\PSD(d)$, with $G$ acting by
$A \cdot \Sigma$, is itself an axiom: an $A \in G \cap \Orth(d)$ fixes $I$, so (K3) makes $\mu_I$
invariant under $G \cap \Orth(d)$, and when $G \supseteq \Orth(d)$ the family has a
rotation-invariant member. \ref{lem:parameter-space}(3) shows that this is exactly what indexing an
orbit of $\GL(d)$ by $\PD(d)$ requires, and (4) there gives an orbit with no isotropic member, which
has no such index. Nondegeneracy is asked at the one pair $(0, I)$; it makes the stabilizer of $\mu_I$
compact when $G \supseteq \Orth(d)$ (for $d \ge 2$, $G \supseteq \mathrm{SO}(d)$ suffices), and need not
otherwise (\ref{lem:parameter-space}(5)).
The isotropic corner, read on the slice $\{sI\}$, is \ref{def:isotropic-kernel-family}, which
leaves the gauge $S_\lambda$ free and asks (ND) for every $s < t$, up to a change of the scale
coordinate unique up to a constant (\ref{prop:isotropic-corner}). \textbf{Not claimed}: any
statement about the one-parameter orders $e^{tB}\Sigma_0e^{tB^{\tp}}$ of the plan's table, which
are not orders on all of $\PSD(d)$ and are not corners of this definition.)}
\end{definition}

% CHANGED Q-0281: split from def:gw-cascade-family (ADR-0025; records/formalization/PRICING.md,
% finding 4). The sentences the definition displays as facts, stated as a lemma: a Lean definition
% cannot carry a claim. Clause (1) is the definition's "Since W contains the rays, ..."; clause (2) its
% "The Loewner order ... and the ray order ... contain the rays and are invariant under every subgroup
% of GL(d)", read with what calling them orders and corners asserts (partial orders on PSD(d); with a
% closed subgroup containing the dilations, wedge data); clause (3) its transform identity, which
% def:affine-cascade-family and def:isotropic-kernel-family display as well. No new mathematics.
\begin{lemma}[the claims of the $(G, W)$ scheme]
\label{lem:gw-scheme-claims}
\uses{def:gw-cascade-family}
\lean{AffineHemigroup.wedge_zero_le, AffineHemigroup.wedge_zero_one_mem,
AffineHemigroup.isWedgeDatum_loewner, AffineHemigroup.isWedgeDatum_ray,
AffineHemigroup.glSet_isClosedSubgroupGL, AffineHemigroup.charFun_push}\leanok
Let $d \ge 1$.
\begin{enumerate}
\item If $(G, W)$ is a wedge datum in dimension $d$, then $0 \le_W \Sigma$ for every
$\Sigma \in \PSD(d)$; in particular $(0, I) \in P_W$.
\item The Loewner order and the ray order are partial orders on $\PSD(d)$ that contain the rays
and are invariant under every $A \in \GL(d)$. Consequently, for every closed subgroup $G$ of
$\GL(d)$ containing the dilations, $(G, \text{Loewner})$ and $(G, \text{rays})$ are wedge data; in
particular $(\GL(d), \text{Loewner})$ and $(\GL(d), \text{rays})$ are.
\item For every probability measure $\mu$ on $\RR^d$, every $A \in \GL(d)$ and every
$\omega \in \RR^d$, $\widehat{A_*\mu}(\omega) = \widehat\mu(A^{\tp}\omega)$.
\end{enumerate}
\statusT\quad\emph{(Cites nothing. Split from \ref{def:gw-cascade-family} on 2026-10-04 (Q-0281,
ADR-0025): the definition displays these sentences, and a definition in Lean cannot carry a claim,
so they are a node of their own; the definition's text is unchanged. Clause (3) is also the identity
\ref{def:affine-cascade-family} and \ref{def:isotropic-kernel-family} display. \textbf{Not
claimed}: that the similarities form a closed subgroup of $\GL(d)$, which the corners
$(\mathrm S(d), \text{rays})$ and $(\mathrm S(d), \text{Loewner})$ need and
\ref{prop:similarity-ray-families}(1) proves.)}
\end{lemma}

\begin{proof}
\leanok
(1) $W$ contains the rays: with $c = 0$ and $c' = 1$, $0 = 0\,\Sigma \le_W 1\,\Sigma = \Sigma$. For
$\Sigma = I \in \PSD(d)$ this is $(0, I) \in P_W$.

(2) \emph{The Loewner order.} It is reflexive, as $0 \in \PSD(d)$; antisymmetric, as a symmetric
$X$ with $X$ and $-X$ positive semidefinite has $\langle x, Xx\rangle = 0$ for every $x$ and so is
$0$; and transitive, as a sum of positive semidefinite matrices is positive semidefinite. For
$A \in \GL(d)$, $A\Sigma'A^{\tp} - A\Sigma A^{\tp} = A(\Sigma' - \Sigma)A^{\tp}$ is positive
semidefinite with $\Sigma' - \Sigma$, since $\langle x, AXA^{\tp}x\rangle = \langle A^{\tp}x,
XA^{\tp}x\rangle$. For $0 \le c \le c'$, $c'\Sigma - c\Sigma = (c' - c)\Sigma$ is positive
semidefinite. \emph{The ray order.} It is reflexive ($c = 1$). Transitive: if $\Sigma = 0$ there is
nothing to show; otherwise $\Sigma' = c\Sigma \ne 0$, so $\Sigma'' = c'\Sigma' = cc'\Sigma$ with
$cc' \ge 1$. Antisymmetric: if $\Sigma \le_W \Sigma'$ and $\Sigma' \le_W \Sigma$ and one of them is
$0$, the other relation forces the second to be $0$ as well (a multiple of $0$ is $0$); otherwise
$\Sigma' = c\Sigma$ and $\Sigma = c'\Sigma'$ give $\Sigma = cc'\Sigma$ with $\Sigma \ne 0$, so
$cc' = 1$, hence $c = c' = 1$ and $\Sigma' = \Sigma$. Invariant: $A \cdot 0 = 0$ and
$A \cdot (c\Sigma) = c\,(A \cdot \Sigma)$. Rays: for $0 \le c \le c'$, $c\Sigma = 0$ if $c = 0$,
and otherwise $c'\Sigma = (c'/c)(c\Sigma)$ with $c'/c \ge 1$. \emph{Wedge data.} Both orders are
invariant under every element of $\GL(d)$, hence of $G$. $\GL(d)$ is closed in itself and contains
$\lambda I$ for $\lambda > 0$.

(3) By the change of variables $y = Ax$,
$\widehat{A_*\mu}(\omega) = \int e^{i\langle\omega,Ax\rangle}\mu(dx)
= \int e^{i\langle A^{\tp}\omega,x\rangle}\mu(dx) = \widehat\mu(A^{\tp}\omega)$.
\end{proof}

% draft: section 2.2, Definition 2.4, restated on kernels. It is the kernel form of
% def:gw-cascade-family at (GL(d), Loewner), written out because every node below reads (K1)-(K4)
% in this form; the passage from operators is lem:operators-to-kernels(4).
% CHANGED Q-0281: typed (the statement skeleton). The displayed claims -- the action preserves the
% order, the transform of an image measure, eq:affine-gaussian and the support of the degenerate law --
% are lem:gw-scheme-claims(2)-(3) and the split lem:affine-gaussian-transform (ADR-0025; PRICING
% finding 4); the text is unchanged and the \lean tag names their declarations too.
% CHANGED Q-0281: blind review R1-5: the tag names glSet_isClosedSubgroupGL, which isWedgeDatum_loewner
% needs at GL(d) for "which preserves the order".
% CHANGED Q-0281 wave 2: blind review W2C-1: the annotation's "that theorem is not used" became false
% when the proof of thm:cascade-infinitely-divisible(1) took the Lean route through Levy's continuity
% theorem; the author's decision of 2026-10-05 accepts that route (Mathlib, no ledger entry), and the
% annotation says so. The statement text is unchanged.
% CHANGED Q-0469 (wording only; records/affine/GATE-affine.md, Q-0317 row "§3 24"): the preamble no
% longer repeats def:gw-cascade-family's (PSD(d), PD(d), the Loewner order, the action, the pushforward,
% the transform and its identity under the action); it points there. "Which preserves the order" and
% the transform identity remain claims of lem:gw-scheme-claims(2)-(3), and the \lean tag is kept as it
% was. No claim and no definition changes.
\begin{definition}[affine cascade family]
\label{def:affine-cascade-family}
\uses{def:gw-cascade-family}
\lean{AffineHemigroup.IsAffineCascadeFamily, AffineHemigroup.affineGaussian,
AffineHemigroup.isWedgeDatum_loewner, AffineHemigroup.glSet_isClosedSubgroupGL,
AffineHemigroup.charFun_push,
AffineHemigroup.charFun_affineGaussian,
AffineHemigroup.affineGaussian_carried_by_range}\leanok
Fix an integer $d \ge 1$, and let $\PSD(d)$, $\PD(d)$, the Loewner order $\Sigma \le \Sigma'$, the
action $A \cdot \Sigma = A\Sigma A^{\tp}$ of $\GL(d)$, the pushforward $A_*\mu$ and the transform
$\widehat\mu$ be as in \ref{def:gw-cascade-family}. An \emph{affine cascade family in dimension
$d$} is a $(\GL(d), \text{Loewner})$-cascade family of kernels (\ref{def:gw-cascade-family}(2)),
that is, a family $(\mu_{\Sigma,\Sigma'})$ of probability measures on $\RR^d$, indexed by the pairs
$\Sigma \le \Sigma'$ in $\PSD(d)$, such that
\begin{enumerate}
\item[(K1)] \emph{(cascade)} $\mu_{\Sigma,\Sigma} = \delta_0$, and
$\mu_{\Sigma,\Sigma''} = \mu_{\Sigma',\Sigma''} * \mu_{\Sigma,\Sigma'}$ whenever
$\Sigma \le \Sigma' \le \Sigma''$;
\item[(K2)] \emph{(continuity)} for every $\omega \in \RR^d$ the map
$(\Sigma,\Sigma') \mapsto \widehat\mu_{\Sigma,\Sigma'}(\omega)$ is continuous on the set of
pairs $\Sigma \le \Sigma'$;
\item[(K3)] \emph{(affine covariance)} $A_*\mu_{\Sigma,\Sigma'} = \mu_{A\Sigma A^{\tp},\,A\Sigma'A^{\tp}}$
for every $A \in \GL(d)$ and every pair;
\item[(K4)] \emph{(nondegeneracy)} $\mu_{0,I} \ne \delta_0$.
\end{enumerate}
We write $\mu_\Sigma := \mu_{0,\Sigma}$ for the \emph{kernel from scale zero}. For $b > 0$ the
\emph{affine Gaussian scale space with constant $b$} is the family in which $\mu_{\Sigma,\Sigma'}$
is the centred Gaussian law with covariance matrix $2b(\Sigma' - \Sigma)$, that is
\begin{equation}\label{eq:affine-gaussian}
  \widehat\mu_{\Sigma,\Sigma'}(\omega) = e^{-b\langle\omega,(\Sigma'-\Sigma)\omega\rangle} ,
\end{equation}
the law being degenerate, carried by the range of $\Sigma' - \Sigma$, when that matrix is singular.
\statusT\quad\emph{(Vocabulary, and no claim. The draft's Definition~2.4 states the axioms on
operators $\Phi_{\Sigma,\Sigma'}$ on $L^1(\RR^d)$: (A1), (A2) for $a \in \RR^d$, (A4) and (A5) of
[V, Def.~3.1], the Loewner cascade, continuity in $L^1$, nondegeneracy, and
$D_A\Phi_{\Sigma,\Sigma'}D_A^{-1} = \Phi_{A\Sigma A^{\tp},A\Sigma'A^{\tp}}$ with
$D_Af = |\det A|^{-1}f(A^{-1}\,\cdot)$; that is the operator form of \ref{def:gw-cascade-family} at
$(\GL(d), \text{Loewner})$, and this node is the kernel form there, written out. It is the form every
proof below reads, and the
passage from the operator form to it is \ref{lem:operators-to-kernels}(4): each operator is
convolution by a unique probability measure, and those measures satisfy (K1)--(K4). That lemma
proves [V, Lem.~4.1] in $\RR^d$ rather than citing it, the released lemma being stated on the
line. The way back, from (K2) to continuity in $L^1$, is \ref{lem:kernels-to-operators}(2), and
\ref{thm:affine-rigidity-operators} reads it.
The parameter set is $\PSD(d)$ with the
origin included, which is the choice the draft's check item~4 proposes: it makes the kernels from scale zero members of the family. (K2) asks for continuity of the transforms pointwise, which is
weak continuity by \Levy's continuity theorem on $\RR^d$. The definition does not need that reading;
the one proof that uses the theorem, that of \ref{thm:cascade-infinitely-divisible}(1), takes it,
with the portmanteau inequality, as a standard theorem machine-checked in Mathlib, not as a ledger
entry (the author's decision of 2026-10-05). Positivity is the word
``probability'' in the first line; the draft's Remark~2.12 is the reason it is load-bearing.)}
\end{definition}

% CHANGED Q-0281: split from def:affine-cascade-family (ADR-0025; records/formalization/PRICING.md,
% finding 4): the definition's claims about its Gaussian member, eq:affine-gaussian ("that is") and the
% sentence on the degenerate law, stated as a lemma. The other two claims the definition displays (the
% action preserves the order; the transform of an image) are lem:gw-scheme-claims(2)-(3).
\begin{lemma}[the affine Gaussian scale space: transform and support]
\label{lem:affine-gaussian-transform}
\uses{def:affine-cascade-family}
\lean{AffineHemigroup.charFun_affineGaussian,
AffineHemigroup.affineGaussian_carried_by_range}\leanok
Let $d \ge 1$, $b > 0$, and $\Sigma \le \Sigma'$ in $\PSD(d)$. The centred Gaussian law with
covariance matrix $2b(\Sigma' - \Sigma)$ has the transform \ref{eq:affine-gaussian}, and it is
carried by the range of $\Sigma' - \Sigma$, which is a proper subspace when that matrix is
singular.
\statusT\quad\emph{(Cites nothing beyond the transform of a centred Gaussian law. Split from
\ref{def:affine-cascade-family} on 2026-10-04 (Q-0281, ADR-0025); the definition's text is
unchanged.)}
\end{lemma}

\begin{proof}
\leanok
Put $X := 2b(\Sigma' - \Sigma)$, positive semidefinite, and let $Z$ be standard normal on $\RR^d$.
The law in question is the image of the law of $Z$, whose transform is $e^{-|\xi|^2/2}$, under
$X^{1/2}$; by the transform of an image measure its transform is
$e^{-|X^{1/2}\omega|^2/2} = e^{-\langle\omega,X\omega\rangle/2}
= e^{-b\langle\omega,(\Sigma'-\Sigma)\omega\rangle}$, which is \ref{eq:affine-gaussian}. The values of $X^{1/2}Z$ lie in the range of $X^{1/2}$, which is the
range of $X$ (the two symmetric matrices have the same kernel), and that is the range of
$\Sigma' - \Sigma$; it is a proper subspace exactly when $\Sigma' - \Sigma$ is singular.
\end{proof}

% draft: the passage from the operator axioms of Definition 2.4 (section 2.2) and of Definition 2.1
% (section 2.1) to kernels, which the draft takes from [V, Lem. 4.1] "read in R^d". Provenance:
% Q-0207, on the author's answer (c) to Q-0205 (2026-09-28). Proved, not cited, and no ledger
% entry: notes/PLAN-d-dimensional-port.md S1 checked each step of [V, Lem. 4.1]'s proof in R^d,
% and the only one-dimensional devices were the interval approximants and their mass-moving bound,
% which become balls here. Clause (3) is the first step of [V, Lem. 4.2]'s proof (the Gaussian
% test function); clause (2) is [V, Lem. 6.1] (1)<=>(2) with GL(d) in place of the dilations. The
% lemma is per operator in (1)-(3), so (4) and (5) are two readings of it: one node serves both
% definitions. It sits after def:affine-cascade-family because it \uses both kernel definitions;
% neither definition \uses it, so the graph stays acyclic. Q-0265 (plan item D2) generalized (4)
% from (GL(d), Loewner) to every wedge datum (G, W) of def:gw-cascade-family; the proof is the same,
% with "A in G" for "A in GL(d)" and the G-invariance of W for that of the Loewner order.
% CHANGED Q-0281: typed (the statement skeleton). Clauses (4)-(5), the assembly into kernel families,
% are split to lem:operators-to-kernels-families (ADR-0025; records/formalization/PRICING.md: the
% per-operator analysis of (1)-(3) against the assembly); the text is unchanged and the \lean tag
% names the child's declarations too. The d-dimensional L^1 operators T_a and D_A are not built: the
% Lean states (A2) and covariance relationally (AffineHemigroup.Defs.Operators), and "the measures that
% (1) gives" as any measures representing the operators, which (1)'s uniqueness makes the same.
% CHANGED Q-0281 wave M, phase 1: the sentence above about T_a and D_A is superseded. They are built
% (the shared library ScaleSpaceCore v0.3.0), and the declarations are re-typed on them, with the
% same mathematics: Phi f = mu * f is the operator equation Phi = ScaleSpace.mconvL1E mu, (A2) is
% Phi T_a = T_a Phi, (2) is D_A(mu * f) = (A_* mu) * D_A f and D_A (mu * .) D_A^{-1} = nu * . as
% written, and mu * gamma is the L^1 element ScaleSpace.mconvL1E mu (gaussL1 d). The text is
% unchanged; records/formalization/SKELETON.md, "Wave M, phase 1", lists every changed type.
% CHANGED Q-0281 wave M, phase 2: proved (AffineHemigroup.OperatorsKernels, the representation
% through AffineHemigroup.RepresentationSpace, the d-dimensional port of [V, Lem. 4.1]'s Lean with
% balls for intervals; (4)-(5) in AffineHemigroup.OperatorFamilies), Lean core. The Lean follows the
% proof of record; two steps are presented differently and need no change of text: the
% identification tests Phi f - mu * f against the indicator of every set of finite measure rather
% than its sign, and (K2) in (4) is read without sequences, through the continuous map
% p |-> (Phi_p gamma)^(omega) / gamma^(omega), which is the argument of (3). The text is unchanged.
% CHANGED Q-0406: \uses gains def:isotropic-operator-family, the axioms of clause (5) as a
% definition node of chapter 3 (the isotropic paper opens with it). The text is unchanged.
% CHANGED Q-0327: notation only (blind draft review flag 32): the standard Gaussian density is
% $\varphi$, no longer $\gamma$, here and in lem:kernels-to-operators.
\begin{lemma}[from operators to kernels in dimension $d$]
\label{lem:operators-to-kernels}
\uses{def:gw-cascade-family,def:affine-cascade-family,def:isotropic-kernel-family,def:isotropic-operator-family}
\lean{AffineHemigroup.operators_to_kernels_repr, AffineHemigroup.operators_to_kernels_unique,
AffineHemigroup.operators_to_kernels_converse, AffineHemigroup.operators_to_kernels_transport,
AffineHemigroup.operators_to_kernels_transport_iff,
AffineHemigroup.operators_to_kernels_continuity, AffineHemigroup.operators_to_kernels_gw,
AffineHemigroup.operators_to_kernels_affine,
AffineHemigroup.operators_to_kernels_isotropic, AffineHemigroup.gaussDensity,
AffineHemigroup.gaussL1}\leanok
Fix $d \ge 1$. On the real space $L^1(\RR^d)$ write $T_af := f(\cdot - a)$ for $a \in \RR^d$,
$D_Af := |\det A|^{-1}f(A^{-1}\,\cdot)$ for $A \in \GL(d)$, and, for a probability measure $\mu$
on $\RR^d$, $(\mu * f)(x) := \int f(x - y)\,\mu(dy)$. Let $\varphi(x) := (2\pi)^{-d/2}e^{-|x|^2/2}$
be the standard Gaussian density.
\begin{enumerate}
\item \emph{(Representation.)} Let $\Phi : L^1(\RR^d) \to L^1(\RR^d)$ be linear and satisfy the
single-operator axioms of \restated{V}{Def.~3.1}{def:line-family} other than the reflection axiom (A3), read in $\RR^d$:
(A1) $\Phi$ is
bounded; (A2) $\Phi T_a = T_a\Phi$ for every $a \in \RR^d$; (A4) $\Phi f \ge 0$ whenever
$f \ge 0$; (A5) $\int\Phi f = \int f$ whenever $f \ge 0$. Then there is a unique probability
measure $\mu$ on $\RR^d$ with
\[
  \Phi f = \mu * f, \qquad f \in L^1(\RR^d) .
\]
Uniqueness holds in the strong form: if $\mu$ and $\nu$ are probability measures on $\RR^d$ with
$\mu * \varphi = \nu * \varphi$, then $\mu = \nu$. Conversely, for every probability measure $\mu$
on $\RR^d$ the operator $f \mapsto \mu * f$ satisfies (A1), (A2), (A4) and (A5).
\item \emph{(Transport.)} For every probability measure $\mu$ on $\RR^d$, every $A \in \GL(d)$ and
every $f \in L^1(\RR^d)$, $D_A(\mu * f) = (A_*\mu) * D_Af$. Consequently, for probability measures
$\mu$ and $\nu$, the operators satisfy $D_A\,(\mu * \cdot\,)\,D_A^{-1} = \nu * \cdot$ if and only
if $A_*\mu = \nu$.
\item \emph{(Continuity.)} If $\mu_n$ and $\mu$ are probability measures on $\RR^d$ with
$\mu_n * \varphi \to \mu * \varphi$ in $L^1(\RR^d)$, then $\widehat\mu_n(\omega) \to
\widehat\mu(\omega)$ for every $\omega \in \RR^d$.
\item \emph{(The $(G, W)$ families, and the affine family.)} Let $(G, W)$ be a wedge datum in
dimension $d$ and $(\Phi_{\Sigma,\Sigma'})_{(\Sigma,\Sigma') \in P_W}$ a $(G, W)$-cascade family of
operators (\ref{def:gw-cascade-family}(1)). Then the measures $\mu_{\Sigma,\Sigma'}$ that (1) gives,
$\Phi_{\Sigma,\Sigma'}f = \mu_{\Sigma,\Sigma'} * f$, form a $(G, W)$-cascade family of kernels
(\ref{def:gw-cascade-family}(2)). In particular, let $(\Phi_{\Sigma,\Sigma'})$ be linear
operators on $L^1(\RR^d)$, indexed by the pairs $\Sigma \le \Sigma'$ in $\PSD(d)$, with the axioms
of the case $(\GL(d), \text{Loewner})$: each $\Phi_{\Sigma,\Sigma'}$ satisfies (A1), (A2),
(A4) and (A5); (A6) $\Phi_{\Sigma,\Sigma} = \mathrm{Id}$ and
$\Phi_{\Sigma',\Sigma''}\Phi_{\Sigma,\Sigma'} = \Phi_{\Sigma,\Sigma''}$ whenever
$\Sigma \le \Sigma' \le \Sigma''$; (A7) for every $f \in L^1(\RR^d)$ the map
$(\Sigma,\Sigma') \mapsto \Phi_{\Sigma,\Sigma'}f$ is continuous from the set of pairs into
$L^1(\RR^d)$; $D_A\Phi_{\Sigma,\Sigma'}D_A^{-1} = \Phi_{A\Sigma A^{\tp},\,A\Sigma'A^{\tp}}$ for
every $A \in \GL(d)$ and every pair; and $\Phi_{0,I} \ne \mathrm{Id}$. Then the measures form an
affine cascade family in dimension $d$ (\ref{def:affine-cascade-family}).
\item \emph{(The isotropic family.)} Let $(\Phi_{s,t})_{0 \le s \le t}$ be linear operators on
$L^1(\RR^d)$ with these axioms: each $\Phi_{s,t}$ satisfies (A1),
(A2), (A4) and (A5); (A6) $\Phi_{t,t} = \mathrm{Id}$ and $\Phi_{s,t}\Phi_{r,s} = \Phi_{r,t}$ for
$r \le s \le t$; (A7) for every $f \in L^1(\RR^d)$ the map $(s,t) \mapsto \Phi_{s,t}f$ is
continuous from $\{0 \le s \le t\}$ into $L^1(\RR^d)$; $\Phi_{s,t}R_Q = R_Q\Phi_{s,t}$ for every
$Q \in \Orth(d)$, with $R_Qf := f(Q^{-1}\,\cdot)$; (A8) there are increasing bijections
$S_\lambda$ of $[0,\infty)$, $\lambda > 0$, with
$D_\lambda\Phi_{s,t} = \Phi_{S_\lambda s,\,S_\lambda t}D_\lambda$, where
$D_\lambda f := \lambda^{-d}f(\cdot/\lambda)$; and (ND) $\Phi_{s,t} \ne \mathrm{Id}$ whenever
$s < t$. Then the measures $\mu_{s,t}$ that (1) gives form an isotropic cascade family in
dimension $d$ (\ref{def:isotropic-kernel-family}).
\end{enumerate}
\statusT\quad\emph{(Proved, not cited, and cites no source: it reads, as [V, Lem.~4.1] does,
Prokhorov's theorem on $\RR^d$, and, as \ref{lem:radial-marginal} does, the uniqueness of the
Fourier transform of a finite measure on $\RR^d$; neither is a ledger entry, and
\texttt{blueprint/AXIOMS.md} is unchanged. The proof of (1) is [V, Lem.~4.1]'s with balls in
place of intervals, the Wendel step specialised to $\RR^d$ and sharpened by (A4) and (A5) to a
probability measure; (3) is the first step of [V, Lem.~4.2]'s proof; (2) is [V, Lem.~6.1]
(1)$\Leftrightarrow$(2), with $\GL(d)$ in place of the dilations. Clauses (1)--(3) concern one
operator and never the parameter set, which is why one node serves every wedge datum of
\ref{def:gw-cascade-family}, the affine family indexed by Loewner pairs among them, and the isotropic
one, indexed by scales. \textbf{Readings.} The draft's
Definition~2.4 leaves its parameter set to its check item~4; (4) takes $\PSD(d)$ with the origin,
as \ref{def:affine-cascade-family} does, and it reads the draft's unspecified
``nondegeneracy'' as $\Phi_{0,I} \ne \mathrm{Id}$, the operator form of (K4). Rotation invariance
in (5) is covariance under $D_Q = R_Q$, since $|\det Q| = 1$, and $D_\lambda = D_{\lambda I}$.
\textbf{Not claimed here}: the converses of (4) and (5), that a kernel family defines by convolution
an operator family with the draft's axioms. They are \ref{lem:kernels-to-operators}(2) and (3),
whose continuity in $L^1$ comes from Gaussian smoothing and Scheff\'e's lemma, without \Levy's
continuity theorem. The sufficiency halves of the module's theorems are stated on kernels, and
\ref{thm:affine-rigidity-operators} reads both directions: this lemma's (1) and (4), and the
converse.)}
\end{lemma}

\begin{proof}
\leanok
Two facts are used throughout. For a probability measure $\mu$ and $f \in L^1(\RR^d)$, $\mu * f$
is the $L^1(\RR^d)$-valued Bochner integral $\int T_yf\,\mu(dy)$: the map $y \mapsto T_yf$ is
continuous because translation acts continuously on $L^1(\RR^d)$, and it is bounded by
$\|f\|_1$. And with $\widehat h(\omega) := \int e^{i\langle\omega,x\rangle}h(x)\,dx$ for
$h \in L^1(\RR^d)$, Fubini gives the pairing
\[
  \widehat{\mu * f}(\omega) = \widehat\mu(\omega)\,\widehat f(\omega),
  \qquad |\widehat h(\omega)| \le \|h\|_1 ,
\]
and $\widehat\varphi(\omega) = e^{-|\omega|^2/2} > 0$.

\emph{(1), the interchange identity.} For $f, g \in L^1(\RR^d)$, $f * g = \int f(y)\,T_yg\,dy$ as
a Bochner integral, dominated by $|f(y)|\,\|g\|_1$. A bounded operator commutes with an
$L^1(\RR^d)$-valued integral, so (A1) and (A2) give $\Phi(f * g) = \int f(y)\,T_y\Phi g\,dy
= f * \Phi g$.

\emph{The approximants.} Let $B_\varepsilon := \{|x| < \varepsilon\}$, $|B_\varepsilon|$ its
volume, $\rho_\varepsilon := |B_\varepsilon|^{-1}\mathbf 1_{B_\varepsilon}$ and
$h_\varepsilon := \Phi\rho_\varepsilon$. By (A4) and (A5), $h_\varepsilon \ge 0$ with
$\int h_\varepsilon = 1$, so $\nu_\varepsilon := h_\varepsilon\,dx$ is a probability measure on
$\RR^d$. For every $f$, $\|f * \rho_\varepsilon - f\|_1 \le \sup_{|y| < \varepsilon}
\|T_yf - f\|_1 \to 0$ as $\varepsilon \downarrow 0$.

\emph{Tightness.} By the interchange identity $\rho_1 * h_\varepsilon = \Phi(\rho_1 *
\rho_\varepsilon)$, which tends to $\Phi\rho_1$ in $L^1(\RR^d)$ by (A1). Along
$\varepsilon_n := 1/(n+1)$ the functions $\rho_1 * h_{\varepsilon_n}$ therefore have uniformly
small tails: given $\eta > 0$ there is $R$ with $\int_{|x| > R}(\rho_1 * h_{\varepsilon_n}) < \eta$
for every $n$ (choose $N$ with $\|\rho_1 * h_{\varepsilon_n} - \Phi\rho_1\|_1 < \eta/2$ for
$n \ge N$, then $R$ for $\Phi\rho_1$ and the first $N$ terms). Convolution with a nonnegative
unit-mass density carried by the unit ball moves mass by at most $1$: for $|y| > R + 1$ the
translate $\rho_1(\cdot - y)$ is carried by $\{|x| > R\}$, so by Tonelli
\[
  \int_{|x| > R}(\rho_1 * h_\varepsilon)(x)\,dx
  = \int h_\varepsilon(y)\int_{|x| > R}\rho_1(x - y)\,dx\,dy
  \ge \nu_\varepsilon\bigl(\{|y| > R + 1\}\bigr) .
\]
So the sequence $(\nu_{\varepsilon_n})$ is tight, with the closed balls of radius $R + 1$ as
compacts.

\emph{The limit.} By Prokhorov's theorem a subsequence $\nu_{\varepsilon_{n_k}}$ converges weakly
to a probability measure $\mu$ on $\RR^d$.

\emph{Identification.} Let $\Psi \in L^\infty(\RR^d)$, acting on $L^1(\RR^d)$ by
$\Psi(h) := \int\Psi h$, and $f \in L^1(\RR^d)$. The map $y \mapsto \Psi(T_yf)$ is continuous and
bounded by $\|\Psi\|_\infty\|f\|_1$, hence a test function for weak convergence. Since
$f * h_\varepsilon = \int h_\varepsilon(y)\,T_yf\,dy$ and $\Psi$ commutes with the integral,
\[
  \Psi(f * h_{\varepsilon_{n_k}}) = \int\Psi(T_yf)\,\nu_{\varepsilon_{n_k}}(dy)
  \ \xrightarrow{\ k \to \infty\ }\ \int\Psi(T_yf)\,\mu(dy) = \Psi(\mu * f) .
\]
On the other hand $f * h_\varepsilon = \Phi(f * \rho_\varepsilon) \to \Phi f$ in $L^1(\RR^d)$, by
the interchange identity and (A1). Hence $\Psi(\Phi f - \mu * f) = 0$ for every $\Psi$; taking
$\Psi$ the sign of $\Phi f - \mu * f$ gives $\Phi f = \mu * f$.

\emph{Uniqueness.} If $\mu * \varphi = \nu * \varphi$, the pairing gives
$\widehat\mu(\omega)\,\widehat\varphi(\omega) = \widehat\nu(\omega)\,\widehat\varphi(\omega)$, so
$\widehat\mu = \widehat\nu$ because $\widehat\varphi > 0$, and $\mu = \nu$ because a finite measure
on $\RR^d$ is determined by its Fourier transform. In particular two measures representing the
same operator agree.

\emph{Converse.} For a probability measure $\mu$, $\|\mu * f\|_1 \le \|f\|_1$ by Minkowski's
integral inequality (A1); $T_a(\mu * f) = \mu * T_af$, both being
$x \mapsto \int f(x - a - y)\,\mu(dy)$ (A2); $\mu * f \ge 0$ when $f \ge 0$ (A4); and
$\int\mu * f = \int f$ for $f \ge 0$ by Tonelli (A5).

\emph{(2).} For $x \in \RR^d$, substituting $z = Ay$,
\[
  \bigl((A_*\mu) * D_Af\bigr)(x) = \int |\det A|^{-1}f\bigl(A^{-1}(x - z)\bigr)\,(A_*\mu)(dz)
  = |\det A|^{-1}\int f(A^{-1}x - y)\,\mu(dy) = D_A(\mu * f)(x) .
\]
$D_A$ is invertible on $L^1(\RR^d)$ with inverse $D_{A^{-1}}$, so applying the identity to
$D_A^{-1}g$ gives $D_A\,(\mu * D_A^{-1}g) = (A_*\mu) * g$ for every $g$: the operator
$D_A\,(\mu * \cdot\,)\,D_A^{-1}$ is convolution by $A_*\mu$. It is convolution by $\nu$ if and
only if $A_*\mu = \nu$, by the uniqueness clause of (1).

\emph{(3).} By the pairing, $|\widehat\mu_n(\omega) - \widehat\mu(\omega)|\,\widehat\varphi(\omega)
\le \|\mu_n * \varphi - \mu * \varphi\|_1 \to 0$; divide by $\widehat\varphi(\omega) > 0$. The test
density must have a zero-free transform, which is why the indicator of a ball would not do.

\emph{(4).} By (1), each $\Phi_{\Sigma,\Sigma'}$ is convolution by a unique probability measure
$\mu_{\Sigma,\Sigma'}$. The four axioms of the kernel form follow.
\begin{itemize}
\item (K1): $\Phi_{\Sigma,\Sigma} = \mathrm{Id}$ is convolution by $\delta_0$, so
$\mu_{\Sigma,\Sigma} = \delta_0$ by uniqueness. By Fubini
$\mu_{\Sigma',\Sigma''} * (\mu_{\Sigma,\Sigma'} * f) = (\mu_{\Sigma',\Sigma''} *
\mu_{\Sigma,\Sigma'}) * f$, so the cascade says that $\mu_{\Sigma,\Sigma''}$ and
$\mu_{\Sigma',\Sigma''} * \mu_{\Sigma,\Sigma'}$ represent the same operator, and they are equal by
uniqueness.
\item (K2): $P_W$ is a subset of a finite-dimensional space, so continuity on it is sequential. If
$(\Sigma_n,\Sigma'_n) \to (\Sigma,\Sigma')$ in $P_W$, continuity at $f = \varphi$ gives
$\mu_{\Sigma_n,\Sigma'_n} * \varphi \to \mu_{\Sigma,\Sigma'} * \varphi$ in $L^1(\RR^d)$, and (3)
gives $\widehat\mu_{\Sigma_n,\Sigma'_n}(\omega) \to \widehat\mu_{\Sigma,\Sigma'}(\omega)$.
\item (K3): for $A \in G$, $A\cdot\Sigma \le_W A\cdot\Sigma'$ because $W$ is $G$-invariant, so the
right-hand side of the covariance axiom is a member of the family, and (2) with
$\mu = \mu_{\Sigma,\Sigma'}$ and $\nu = \mu_{A\cdot\Sigma,A\cdot\Sigma'}$ turns the axiom into
$A_*\mu_{\Sigma,\Sigma'} = \mu_{A\cdot\Sigma,\,A\cdot\Sigma'}$.
\item (K4): if $\mu_{0,I} = \delta_0$ then $\Phi_{0,I}f = \delta_0 * f = f$ for every $f$, against
$\Phi_{0,I} \ne \mathrm{Id}$.
\end{itemize}
The pair $(\GL(d), \text{Loewner})$ is a wedge
datum: $\GL(d)$ is closed in itself and contains the dilations, and the Loewner order is
$\GL(d)$-invariant, since $A(\Sigma' - \Sigma)A^{\tp}$ is positive semidefinite with
$\Sigma' - \Sigma$, and contains the rays, since $c'\Sigma - c\Sigma = (c' - c)\Sigma$; its kernel
form is \ref{def:affine-cascade-family}.

\emph{(5).} The same argument on the parameter set $\{0 \le s \le t\}$ gives (I1) from (A6) and
(I2) from (A7). (I3): $|\det Q| = 1$, so $R_Q = D_Q$, and the commutation reads
$D_Q\Phi_{s,t}D_Q^{-1} = \Phi_{s,t}$; (2) with $\nu = \mu = \mu_{s,t}$ gives
$Q_*\mu_{s,t} = \mu_{s,t}$. (I4): $D_\lambda = D_{\lambda I}$, and (A8) reads
$D_{\lambda I}\Phi_{s,t}D_{\lambda I}^{-1} = \Phi_{S_\lambda s,\,S_\lambda t}$; (2) gives
$(\lambda I)_*\mu_{s,t} = \mu_{S_\lambda s,\,S_\lambda t}$. (ND): if $\mu_{s,t} = \delta_0$ with
$s < t$ then $\Phi_{s,t} = \mathrm{Id}$, against (ND).
\end{proof}

% CHANGED Q-0281: split from lem:operators-to-kernels (ADR-0025; records/formalization/PRICING.md):
% clauses (4) and (5), the assembly of the per-operator clauses (1)-(3) into the two kernel
% definitions, copied verbatim. The parent's text and numbering are unchanged. Re-copied after the
% merge of main (Q-0287/Q-0289 reworded the parent's clauses without changing what they claim).
% CHANGED Q-0406: \uses gains def:isotropic-operator-family, whose Lean structure the tag below has
% named since the node was typed. The text is unchanged.
\begin{lemma}[from operator families to kernel families]
\label{lem:operators-to-kernels-families}
\uses{lem:operators-to-kernels,def:gw-cascade-family,def:affine-cascade-family,def:isotropic-kernel-family,def:isotropic-operator-family,lem:gw-scheme-claims}
\lean{AffineHemigroup.operators_to_kernels_gw, AffineHemigroup.operators_to_kernels_affine,
AffineHemigroup.operators_to_kernels_isotropic, AffineHemigroup.IsIsotropicOperatorFamily,
AffineHemigroup.IsConvOnPairs}\leanok
Fix $d \ge 1$, and use the notation of \ref{lem:operators-to-kernels}.
\begin{enumerate}
\item[(4)] \emph{(The $(G, W)$ families, and the affine family.)} Let $(G, W)$ be a wedge datum in
dimension $d$ and $(\Phi_{\Sigma,\Sigma'})_{(\Sigma,\Sigma') \in P_W}$ a $(G, W)$-cascade family of
operators (\ref{def:gw-cascade-family}(1)). Then the measures $\mu_{\Sigma,\Sigma'}$ that
\ref{lem:operators-to-kernels}(1) gives, $\Phi_{\Sigma,\Sigma'}f = \mu_{\Sigma,\Sigma'} * f$, form a
$(G, W)$-cascade family of kernels (\ref{def:gw-cascade-family}(2)). In particular, let
$(\Phi_{\Sigma,\Sigma'})$ be linear operators on $L^1(\RR^d)$, indexed by the pairs
$\Sigma \le \Sigma'$ in $\PSD(d)$, with the axioms of the case $(\GL(d), \text{Loewner})$: each
$\Phi_{\Sigma,\Sigma'}$ satisfies (A1), (A2), (A4) and (A5); (A6) $\Phi_{\Sigma,\Sigma} = \mathrm{Id}$ and
$\Phi_{\Sigma',\Sigma''}\Phi_{\Sigma,\Sigma'} = \Phi_{\Sigma,\Sigma''}$ whenever
$\Sigma \le \Sigma' \le \Sigma''$; (A7) for every $f \in L^1(\RR^d)$ the map
$(\Sigma,\Sigma') \mapsto \Phi_{\Sigma,\Sigma'}f$ is continuous from the set of pairs into
$L^1(\RR^d)$; $D_A\Phi_{\Sigma,\Sigma'}D_A^{-1} = \Phi_{A\Sigma A^{\tp},\,A\Sigma'A^{\tp}}$ for
every $A \in \GL(d)$ and every pair; and $\Phi_{0,I} \ne \mathrm{Id}$. Then the measures form an
affine cascade family in dimension $d$ (\ref{def:affine-cascade-family}).
\item[(5)] \emph{(The isotropic family.)} Let $(\Phi_{s,t})_{0 \le s \le t}$ be linear operators on
$L^1(\RR^d)$ with these axioms: each $\Phi_{s,t}$ satisfies (A1),
(A2), (A4) and (A5); (A6) $\Phi_{t,t} = \mathrm{Id}$ and $\Phi_{s,t}\Phi_{r,s} = \Phi_{r,t}$ for
$r \le s \le t$; (A7) for every $f \in L^1(\RR^d)$ the map $(s,t) \mapsto \Phi_{s,t}f$ is
continuous from $\{0 \le s \le t\}$ into $L^1(\RR^d)$; $\Phi_{s,t}R_Q = R_Q\Phi_{s,t}$ for every
$Q \in \Orth(d)$, with $R_Qf := f(Q^{-1}\,\cdot)$; (A8) there are increasing bijections
$S_\lambda$ of $[0,\infty)$, $\lambda > 0$, with
$D_\lambda\Phi_{s,t} = \Phi_{S_\lambda s,\,S_\lambda t}D_\lambda$, where
$D_\lambda f := \lambda^{-d}f(\cdot/\lambda)$; and (ND) $\Phi_{s,t} \ne \mathrm{Id}$ whenever
$s < t$. Then the measures $\mu_{s,t}$ that \ref{lem:operators-to-kernels}(1) gives form an
isotropic cascade family in dimension $d$ (\ref{def:isotropic-kernel-family}).
\end{enumerate}
\statusT\quad\emph{(Cites nothing. Clauses (4) and (5) of \ref{lem:operators-to-kernels}, split off
on 2026-10-04 (Q-0281, ADR-0025): they assemble the per-operator clauses (1)--(3) into the two
kernel definitions, a different kind and cost of argument; the numbering stays the parent's. The Lean
quantifies over any measures representing the operators, which by the uniqueness clause of (1) are
the measures (1) gives.)}
\end{lemma}

\begin{proof}
\leanok
This is the part of the proof of \ref{lem:operators-to-kernels} headed (4) and (5); that $(\GL(d),
\text{Loewner})$ is a wedge datum is \ref{lem:gw-scheme-claims}(2).
\end{proof}

% not in the draft: queue item Q-0314 (hub), brief item 2, on review flag 13
% (records/affine/REVIEW-paper-affine-fixes.md): the converses of lem:operators-to-kernels(4) and (5),
% which that lemma's annotation and def:affine-cascade-family's leave unclaimed because continuity in
% L^1 "from (K2) or (I2) would need Levy's continuity theorem". The brief admitted that theorem as a
% ledger entry A7 for the purpose. What writing the statement down found: the obligation is not weak
% convergence but L^1 convergence of mu_n * f, and Gaussian smoothing gives that directly -- the
% transform of mu_n * gamma_sigma is integrable, so mu_n * gamma_sigma converges pointwise by its
% Fourier integral and dominated convergence, then in L^1 by Scheffe's argument, and
% gamma_sigma * f -> f closes. So the node cites nothing and A7 is not admitted
% (records/affine/READING-sato1999levy-continuity.md keeps the transcription of Sato's Prop. 2.5(vii),
% p. 9, should a later proof need weak convergence itself; Mathlib proves that theorem as
% MeasureTheory.ProbabilityMeasure.tendsto_iff_tendsto_charFun). No existing statement changes: the
% "Not claimed" sentences of lem:operators-to-kernels and def:affine-cascade-family stay true of those
% nodes, and their now stale reason is owed in records/affine/GATE-affine.md. No \lean tag: the Lean
% skeleton session types it.
% Cost, priced when the statement was written: (1) is the only analysis -- a Fourier integral of a
% measure against a Gaussian (Mathlib has the Gaussian's Fourier transform on inner product spaces),
% Fubini, dominated convergence twice, and the approximate identity in L^1 (continuity of
% translation); about 150-250 Lean lines on top of whatever the skeleton provides for mu * f on
% L^1(R^d), which lem:operators-to-kernels encodes relationally. Alternatively (1) follows from
% Mathlib's Levy theorem and the same Scheffe step on compactly supported approximants, at a similar
% cost. (2) and (3) are assembly from lem:operators-to-kernels(1)-(2), about 60-100 lines each once
% mu * f is a declaration.
% CHANGED Q-0281 wave 2: blind review W2C-1: the annotation adds where the module does use Levy's
% continuity theorem (thm:cascade-infinitely-divisible(1), from Mathlib; the author's decision of
% 2026-10-05). The statement text and this node's route are unchanged.
% CHANGED Q-0281 wave M, phase 1: typed (the "No \lean tag" above is superseded). Four declarations
% in AffineSkeleton/Ch04Classical.lean cover the node (ADR-0025): kernels_to_operators_continuity
% is (1); kernels_to_operators_gw is (2) up to its last sentence, the operators being any family
% that is convolution by the kernels on P_W (IsConvOnPairs; such a family exists, and the operator
% axioms read it on P_W only); kernels_operators_correspondence is that last sentence, stated as
% round trips on P_W (the blind referee's note 1); kernels_to_operators_isotropic is (3). The
% operators are the shared library's ScaleSpace.mconvL1E, so the "mu * f is a declaration" the cost
% line above waits for exists. Priced at typing: (1) about 150-250 lines; (2) about 60-100 and the
% correspondence about 40 more; (3) about 60-100, assembly. The text is unchanged.
% CHANGED Q-0281 wave M, phase 2: proved, Lean core, on the proof of record's route, without
% Levy's continuity theorem (AffineHemigroup.KernelsContinuity for (1): the Fourier integral of
% gamma_sigma from Mathlib's Gaussian integral, Fubini, dominated convergence, Scheffe's argument;
% AffineHemigroup.KernelsOperators for (2), (3) and the correspondence, continuity on P_W read
% sequentially through (1)). Two of the three facts are checked by other means than Fubini, which
% the text's "each by Fubini's theorem" leaves open: ||g * f||_1 <= ||g||_1 ||f||_1 bounds the
% L^1-valued integral g * f = int g(y) T_y f dy, and nu * (g * f) = (nu * g) * f follows from the
% commutativity of convolution and the interchange identity. The text is unchanged.
% CHANGED Q-0406: \uses gains def:isotropic-operator-family, the definition whose axioms clause (3)
% cites through lem:operators-to-kernels(5). The text is unchanged.
\begin{lemma}[from kernels to operators in dimension $d$]
\label{lem:kernels-to-operators}
\uses{lem:operators-to-kernels,lem:operators-to-kernels-families,def:gw-cascade-family,def:isotropic-kernel-family,def:isotropic-operator-family,lem:gw-scheme-claims}
\lean{AffineHemigroup.kernels_to_operators_continuity, AffineHemigroup.kernels_to_operators_gw,
AffineHemigroup.kernels_operators_correspondence,
AffineHemigroup.kernels_to_operators_isotropic}\leanok
Fix $d \ge 1$, and use the notation of \ref{lem:operators-to-kernels}.
\begin{enumerate}
\item \emph{(Continuity.)} If $\mu_n$ and $\mu$ are probability measures on $\RR^d$ with
$\widehat\mu_n(\omega) \to \widehat\mu(\omega)$ for every $\omega \in \RR^d$, then
$\mu_n * f \to \mu * f$ in $L^1(\RR^d)$ for every $f \in L^1(\RR^d)$.
\item \emph{(The $(G, W)$ families.)} Let $(G, W)$ be a wedge datum in dimension $d$ and
$(\mu_{\Sigma,\Sigma'})_{(\Sigma,\Sigma') \in P_W}$ a $(G, W)$-cascade family of kernels
(\ref{def:gw-cascade-family}(2)). Then the operators $\Phi_{\Sigma,\Sigma'}f := \mu_{\Sigma,\Sigma'} * f$
on $L^1(\RR^d)$ form a $(G, W)$-cascade family of operators (\ref{def:gw-cascade-family}(1)), and the
measures that \ref{lem:operators-to-kernels}(1) gives for them are the $\mu_{\Sigma,\Sigma'}$. So, for
every wedge datum, convolution is a one-to-one correspondence between the $(G, W)$-cascade families of
kernels and those of operators, and its inverse is the passage of
\ref{lem:operators-to-kernels}(4).
\item \emph{(The isotropic family.)} Let $(\mu_{s,t})_{0 \le s \le t}$ be an isotropic cascade family
in dimension $d$ (\ref{def:isotropic-kernel-family}). Then the operators $\Phi_{s,t}f := \mu_{s,t} * f$
on $L^1(\RR^d)$ satisfy the axioms of \ref{lem:operators-to-kernels}(5), with the bijections
$S_\lambda$ of (I4) in (A8), and the measures that \ref{lem:operators-to-kernels}(1) gives for them are
the $\mu_{s,t}$.
\end{enumerate}
\statusT\quad\emph{(Proved, not cited, and cites no source: it reads Fubini's theorem, dominated
convergence, the Gaussian integral and the continuity of translation in $L^1(\RR^d)$, and, through
\ref{lem:operators-to-kernels}(1)--(2), what that lemma reads; no ledger entry, and
\texttt{blueprint/AXIOMS.md} gains none. \textbf{What writing it down found.} The converses of
\ref{lem:operators-to-kernels}(4)--(5) were left unclaimed because continuity in $L^1$ from
(K2) seemed to need \Levy's continuity theorem on $\RR^d$ (Sato 1999, Prop.~2.5(vii), p.~9, read at
the page in \texttt{records/affine/READING-sato1999levy-continuity.md}). The obligation is weaker than
weak convergence: (1) asks only that $\mu_n * f \to \mu * f$ in $L^1$, and Gaussian smoothing gives
it directly, so this proof does not use the theorem, and no ledger entry is admitted for it. The
module uses it once, in the proof of \ref{thm:cascade-infinitely-divisible}(1), which needs weak
convergence itself; it takes the theorem there from Mathlib, as a standard theorem and not a ledger
entry (the author's decision of 2026-10-05). \textbf{Readings.} The
operators are linear on the real space $L^1(\RR^d)$, as in \ref{def:gw-cascade-family}(1); continuity
on $P_W$ is read sequentially, $P_W$ being a subset of a finite-dimensional space. \textbf{Not
claimed}: weak convergence of the $\mu_n$ in (1), which is not needed; anything about operators on
$L^p(\RR^d)$ for $p > 1$.)}
\end{lemma}

\begin{proof}
\leanok
For $\sigma > 0$ write $\varphi_\sigma(x) := \sigma^{-d}\varphi(x/\sigma)$, the centred Gaussian density
with covariance $\sigma^2I$. For a probability measure $\nu$ and $f, g \in L^1(\RR^d)$ three facts are
used, each by Fubini's theorem: $\|\nu * f\|_1 \le \|f\|_1$; $\|g * f\|_1 \le \|g\|_1\|f\|_1$; and
$\nu * (g * f) = (\nu * g) * f$, where $\nu * g$ is the integrable function
$x \mapsto \int g(x - y)\,\nu(dy)$.

\emph{(1), the smoothed measures.} The Gaussian integral gives, for every $z \in \RR^d$,
\[
  \varphi_\sigma(z) = (2\pi)^{-d}\int e^{-i\langle\omega,z\rangle}e^{-\sigma^2|\omega|^2/2}\,d\omega .
\]
Put $z = x - y$ and integrate in $\nu(dy)$. The integrand is bounded by $e^{-\sigma^2|\omega|^2/2}$,
which is integrable for $d\omega \otimes \nu$, so Fubini's theorem gives, for every probability
measure $\nu$ and every $x \in \RR^d$,
\[
  (\nu * \varphi_\sigma)(x) = (2\pi)^{-d}\int e^{-i\langle\omega,x\rangle}\,\widehat\nu(\omega)\,
  e^{-\sigma^2|\omega|^2/2}\,d\omega .
\]
Put $h_n := \mu_n * \varphi_\sigma$ and $h := \mu * \varphi_\sigma$. Since $|\widehat\mu_n| \le 1$ and
$\widehat\mu_n \to \widehat\mu$ pointwise, dominated convergence gives $h_n(x) \to h(x)$ for every
$x$. The $h_n$ and $h$ are nonnegative with integral $1$ (Tonelli), so $\int(h_n - h) = 0$ and
\[
  \|h_n - h\|_1 = 2\int(h - h_n)^+ \xrightarrow{\ n \to \infty\ } 0 ,
\]
by dominated convergence again, since $0 \le (h - h_n)^+ \le h$ and $(h - h_n)^+ \to 0$ pointwise.

\emph{(1), the approximation.} By Minkowski's integral inequality and the substitution $y = \sigma z$,
\[
  \|\varphi_\sigma * f - f\|_1 \le \int\varphi_\sigma(y)\,\|T_yf - f\|_1\,dy
  = \int\varphi(z)\,\|T_{\sigma z}f - f\|_1\,dz ,
\]
which tends to $0$ as $\sigma \downarrow 0$ by dominated convergence: the integrand is bounded by
$2\|f\|_1\varphi(z)$ and tends to $0$ pointwise, translation being continuous on $L^1(\RR^d)$. Now
split $f = (f - \varphi_\sigma * f) + \varphi_\sigma * f$. By the three facts,
$\mu_n * (\varphi_\sigma * f) = h_n * f$ and $\mu * (\varphi_\sigma * f) = h * f$, so
\[
  \|\mu_n * f - \mu * f\|_1 \le 2\,\|f - \varphi_\sigma * f\|_1 + \|(h_n - h) * f\|_1
  \le 2\,\|f - \varphi_\sigma * f\|_1 + \|h_n - h\|_1\,\|f\|_1 .
\]
Given $\varepsilon > 0$, choose $\sigma$ so that the first term is below $\varepsilon$, and then $n_0$ so
that the second is below $\varepsilon$ for $n \ge n_0$.

\emph{(2).} Each $\Phi_{\Sigma,\Sigma'}$ is linear and satisfies (A1), (A2), (A4) and (A5) by the
converse in \ref{lem:operators-to-kernels}(1). \emph{Cascade (A6).} $\Phi_{\Sigma,\Sigma} = \delta_0 * \cdot
= \mathrm{Id}$ by (K1), and for $\Sigma \le_W \Sigma' \le_W \Sigma''$, by the associativity of
convolution (Fubini) and (K1),
$\Phi_{\Sigma',\Sigma''}\Phi_{\Sigma,\Sigma'}f = \mu_{\Sigma',\Sigma''} * (\mu_{\Sigma,\Sigma'} * f)
= (\mu_{\Sigma',\Sigma''} * \mu_{\Sigma,\Sigma'}) * f = \Phi_{\Sigma,\Sigma''}f$.
\emph{Continuity (A7).} $P_W$ is a subset of the finite-dimensional space of pairs of symmetric matrices,
so continuity on it is sequential. If $(\Sigma_n,\Sigma'_n) \to (\Sigma,\Sigma')$ in $P_W$, (K2) gives
$\widehat\mu_{\Sigma_n,\Sigma'_n}(\omega) \to \widehat\mu_{\Sigma,\Sigma'}(\omega)$ for every $\omega$,
and (1) gives $\Phi_{\Sigma_n,\Sigma'_n}f \to \Phi_{\Sigma,\Sigma'}f$ in $L^1(\RR^d)$ for every $f$.
\emph{(ND).} If $\Phi_{0,I} = \mathrm{Id}$, then $\mu_{0,I}$ and $\delta_0$ represent the same
operator, so $\mu_{0,I} = \delta_0$ by the uniqueness clause of \ref{lem:operators-to-kernels}(1),
against (K4). \emph{Covariance.} For $A \in G$ and $(\Sigma,\Sigma') \in P_W$, the pair
$(A\cdot\Sigma, A\cdot\Sigma')$ lies in $P_W$ because $W$ is $G$-invariant, and (K3) says
$A_*\mu_{\Sigma,\Sigma'} = \mu_{A\cdot\Sigma,\,A\cdot\Sigma'}$; by \ref{lem:operators-to-kernels}(2)
this is $D_A\Phi_{\Sigma,\Sigma'}D_A^{-1} = \Phi_{A\cdot\Sigma,\,A\cdot\Sigma'}$. \emph{The measures.}
$\Phi_{\Sigma,\Sigma'}$ is convolution by $\mu_{\Sigma,\Sigma'}$, so by the uniqueness clause of
\ref{lem:operators-to-kernels}(1) the measure that clause gives is $\mu_{\Sigma,\Sigma'}$. \emph{The
correspondence.} Kernels to operators to kernels returns the kernels, by the last step. Operators to
kernels to operators returns the operators, since the measures of
\ref{lem:operators-to-kernels}(4) satisfy $\Phi_{\Sigma,\Sigma'}f = \mu_{\Sigma,\Sigma'} * f$.

\emph{(3).} The same argument on the parameter set $\{0 \le s \le t\}$, a subset of $\RR^2$, gives the
four single-operator axioms, (A6) from (I1), (A7) from (I2) and (1), and (ND) at every $s < t$ from
(ND). For $Q \in \Orth(d)$, $|\det Q| = 1$, so $D_Q = R_Q$, and (I3) with
\ref{lem:operators-to-kernels}(2) gives $R_Q\Phi_{s,t}R_Q^{-1} = \Phi_{s,t}$, that is
$\Phi_{s,t}R_Q = R_Q\Phi_{s,t}$. For $\lambda > 0$, $D_\lambda = D_{\lambda I}$, and (I4) with
\ref{lem:operators-to-kernels}(2) gives $D_\lambda\Phi_{s,t}D_\lambda^{-1} =
\Phi_{S_\lambda s,\,S_\lambda t}$, which is (A8). The measures are the $\mu_{s,t}$, as in (2).
\end{proof}

% not in the draft: queue item Q-0280 (hub), on the question def:gw-cascade-family left open
% (Q-0265): is the scheme's isotropic corner, read on the slice {sI}, the isotropic theory of
% chapter 3, whose definition leaves the gauge S_lambda free and asks (ND) at every pair? It is, up
% to a change of the scale coordinate that [V, Prop. 6.6] supplies and that is unique up to a
% constant; chapter 3's definition stays as it is, with the reduction to the line. What writing it
% down found: (i) nondegeneracy at the one pair (0, I) gives it at every pair of the slice, so
% chapter 3's (ND) costs nothing there; off the slice it does not follow (the shape-blind family of
% prop:similarity-ray-families(3)(ii) has trivial increments along every singular Sigma, d >= 2);
% (ii) the free gauge is the only other difference, and the variance gauge of [V, Rem. 6.7] removes
% it; (iii) the slice fixes the family for GL(d) and not for the similarities, so chapter 3 is the
% corner's slice, and the corner carries the shapes besides. Placement: here, where the scheme is
% defined. The node reads two nodes of section 5.6 (lem:gw-kernels-from-origin,
% prop:gw-ray-families), forward edges like thm:cascade-infinitely-divisible's, and
% lem:parameter-space(2) of section 4.4; it can move to 5.6, after prop:gw-ray-families, without
% change.
% CHANGED Q-0281 wave 2: blind review W2C-3: the comment above names the nodes the wave-2 proof
% reads (it no longer reads prop:similarity-ray-families); in the proof, that (S(d), rays) is a
% wedge datum moves from (4), after "Let d >= 2", into (2), which needs it and holds for every d >= 1.
% CHANGED Q-0281: blind review R1-3: the tag names isotropic_corner_gauge_forall, (2) for every gauge
% chi with the relation (the proof reads nothing else of chi), beside the existential reading.
% CHANGED Q-0281 wave 2: proved (AffineHemigroup.IsotropicCorner, Lean core), statement unchanged.
% The proof of (4) is rewritten to the Lean route: S(d) is shown closed in GL(d) here, the shape-blind
% family is built through prop:gw-ray-families(2), and the two kernels from scale zero differ because
% Stab(m) = O(d) (lem:parameter-space(2)), so the proof no longer reads prop:similarity-ray-families
% and the \uses edge to it is dropped (the annotation still refers to its family (ii)); the edges to
% lem:gw-scheme-claims and lem:parameter-space are added. (2)'s restriction to S(d) reads only
% S(d) in GL(d).
\begin{proposition}[the slice of the isotropic corner is the isotropic theory]
\label{prop:isotropic-corner}
\uses{def:gw-cascade-family,def:isotropic-kernel-family,lem:isotropic-marginal-family,lem:gw-kernels-from-origin,prop:gw-ray-families,lem:gw-scheme-claims,lem:parameter-space}
\lean{AffineHemigroup.isotropic_corner_slice, AffineHemigroup.isotropic_corner_gauge,
AffineHemigroup.isotropic_corner_gauge_forall, AffineHemigroup.isotropic_corner_gauge_unique, AffineHemigroup.isotropic_corner_rescale,
AffineHemigroup.isotropic_corner_gl_determined,
AffineHemigroup.isotropic_corner_sim_not_determined}\leanok
Let $d \ge 1$, and let $\mathrm S(d) := \{cQ : c > 0,\ Q \in \Orth(d)\}$ be the similarities. Let
$(G, W)$ be a wedge datum in dimension $d$ (\ref{def:gw-cascade-family}) and
$(\nu_{\Sigma,\Sigma'})$ a $(G, W)$-cascade family of kernels, with kernels from scale zero
$\nu_\Sigma := \nu_{0,\Sigma}$. Since $W$ contains the rays, $(sI, tI) \in P_W$ whenever
$0 \le s \le t$; the family $(\nu_{sI,tI})_{0 \le s \le t}$ is the \emph{slice} of
$(\nu_{\Sigma,\Sigma'})$.
\begin{enumerate}
\item \emph{(The slice is an isotropic family.)} If $\Orth(d) \subseteq G$, the slice of every
$(G, W)$-cascade family of kernels is an isotropic cascade family in dimension $d$
(\ref{def:isotropic-kernel-family}) with the dilation action $S_\lambda s = \lambda^2 s$. In
particular $\nu_{sI,tI} \ne \delta_0$ whenever $s < t$.
\item \emph{(Every isotropic family is a slice, in the variance gauge.)} Let $(\mu_{s,t})$ be an
isotropic cascade family in dimension $d$ with dilation action $(S_\lambda)$. There is a gauge
$\chi$ with $\chi(S_\lambda t) = \lambda\,\chi(t)$ for every $\lambda > 0$ and $t \ge 0$; put
$\varphi := \chi^2$, a gauge. Then there is exactly one $(\GL(d), \text{rays})$-cascade family of
kernels $(\nu_{\Sigma,\Sigma'})$ with
\[
  \mu_{s,t} = \nu_{\varphi(s)I,\,\varphi(t)I}, \qquad 0 \le s \le t ,
\]
and its restriction to the datum $(\mathrm S(d), \text{rays})$ is an $(\mathrm S(d), \text{rays})$-cascade family of kernels
with the same property.
\item \emph{(The change of gauge is unique up to a constant.)} In the setting of (2), let $\psi$ be
a gauge, $(G, W)$ a wedge datum, and $(\nu'_{\Sigma,\Sigma'})$ a
$(G, W)$-cascade family of kernels with $\mu_{s,t} = \nu'_{\psi(s)I,\,\psi(t)I}$ for
$0 \le s \le t$. Then $\psi = c\,\varphi$ for some $c > 0$. Conversely, for every $c > 0$ the family
$(\nu_{c^{-1}\Sigma,\,c^{-1}\Sigma'})$ is a $(\GL(d), \text{rays})$-cascade family of kernels with
this property for $\psi = c\,\varphi$.
\item \emph{(What the slice determines.)} A $(\GL(d), \text{rays})$-cascade family of kernels is
determined by its slice. For $d \ge 2$ an $(\mathrm S(d), \text{rays})$-cascade family of kernels
is not: the slice of every such family is the slice of two such families whose kernels from scale zero differ at every $\Sigma \in \PD(d)$ that is not a multiple of $I$.
\end{enumerate}
\statusT\quad\emph{(Cites one released result, [V, Prop.~6.6] (orbit coordinate; the gauge),
restated in the proof where it is used and applied to the marginal family of
\ref{lem:isotropic-marginal-family}. [V] machine-checks it with \texttt{\#print axioms} reducing to
Lean core, so it spends no interface of [V]'s ledger; only its first conclusions are used (the orbit
coordinate and the gauge relation), not its similarity form with the exponent in [V]'s class of
\Levy{} exponents. That a law fixed by one dilation is $\delta_0$ is the argument of [V, Lem.~6.3], there for a
real function, proved here in three lines. Clause (3) holds for every wedge datum: its proof reads
only (K1), (K3) at the dilations, which every wedge datum's group contains, and (K4). It asked
$\Orth(d) \subseteq G$ until the blind check found the hypothesis unused (the author's decision of
2026-10-04 to drop it; Q-0284). No ledger entry of this module. \textbf{What it is worth.} It settles
the question \ref{def:gw-cascade-family} left open: chapter~3's definition is the scheme's isotropic
corner read on the multiples of $I$, and the two ways its axioms differ from the scheme's cost
nothing. Nondegeneracy at the one pair $(0, I)$ gives it at every pair of the slice (clause (1)).
Off the slice it does not follow: for $d \ge 2$ the family (ii) of
\ref{prop:similarity-ray-families}(3) has $\nu_\Sigma = \delta_0$ at every singular $\Sigma$, so its
increments along the ray of a singular $\Sigma$ are all $\delta_0$, and (ND) at every pair would
exclude it; the scheme asks (ND) at $(0, I)$ only. The free gauge $S_\lambda$ of (I4) is the line's, [V]'s (A8). The scheme fixes how
the dilations act on $\PSD(d)$, so that on the slice $\lambda$ acts by $s \mapsto \lambda^2 s$ (the
last sentence of \ref{lem:affine-gauge}), and (2) and (3) say that passing from one to the other is a
change of the scale coordinate, unique up to a constant: the variance gauge $\varphi = \chi^2$ of
[V, Rem.~6.7], $\chi$ being [V]'s canonical gauge of the marginal family. So the article states
chapter~3's definition as it stands and calls it the corner, and the reduction to the line
(\ref{thm:isotropic-classification}) is untouched. Clause (4) says what the corner adds to its
slice: nothing for $\GL(d)$, the shapes for the similarities. With
\ref{thm:isotropic-classification}, clause (3) identifies that theorem's gauge with $\chi$ up to a
constant, since its formula exhibits the square of its gauge as a $\psi$ of (3), for the
$(\GL(d), \text{rays})$-cascade family whose kernel at $I$ has the transform $e^{-F(|\xi|)}$; no node
reads this, and this node does not use the classification. \textbf{Not claimed}: the operator form
(\ref{lem:operators-to-kernels}(4) and (5) take the two operator definitions to their kernel forms,
and the node is stated on kernels); and, beyond clause (3), anything for a group not containing
$\Orth(d)$, whose slice need not satisfy (I3) (for $G$ the dilations and the rays, $\nu_\Sigma := (\Sigma^{1/2})_*m$ with $m$
self-decomposable and not rotation invariant).)}
\end{proposition}

\begin{proof}
\leanok
A fact used twice, the argument of \prog{V}{Lem.~6.3}: \emph{if $m$ is a probability measure on
$\RR^d$ with $(aI)_*m = m$ for some $a > 0$, $a \ne 1$, then $m = \delta_0$.} Indeed
$\widehat m(\xi) = \widehat m(a\xi)$, so $\widehat m(\xi) = \widehat m(a^n\xi)$ for every
$n \in \ZZ$; choosing the sign of $n$ so that $a^n \to 0$, continuity of $\widehat m$ at $0$ gives
$\widehat m(\xi) = 1$, and $m = \delta_0$ by uniqueness of the transform.

\emph{(1).} Let $\Orth(d) \subseteq G$ and put $\mu_{s,t} := \nu_{sI,tI}$. For $0 \le r \le s \le t$,
$rI \le_W sI \le_W tI$ because $W$ contains the rays, so (K1) gives (I1). The map
$(s,t) \mapsto (sI,tI)$ is continuous from $\{0 \le s \le t\}$ into $P_W$, and (K2) composed with it
is (I2). For $Q \in \Orth(d) \subseteq G$, $Q \cdot (sI) = sQQ^{\tp} = sI$, and (K3) gives (I3). For
$\lambda > 0$, $\lambda I \in G$ and $(\lambda I) \cdot (sI) = \lambda^2 s\,I$, so (K3) gives
$(\lambda I)_*\mu_{s,t} = \mu_{\lambda^2 s,\,\lambda^2 t}$, which is (I4) with the increasing
bijection $S_\lambda s := \lambda^2 s$ of $[0,\infty)$. (ND): (K3) at the pair $(0, I)$ with
$A = \sqrt t\,I$ gives $\nu_{tI} = (\sqrt t\,I)_*\nu_I$ for $t > 0$. Let $s < t$ and suppose
$\mu_{s,t} = \delta_0$. By (K1) at $0 \le_W sI \le_W tI$, $\nu_{tI} = \mu_{s,t} * \nu_{sI} = \nu_{sI}$.
If $s = 0$, then $\nu_{0} = \delta_0$ by (K1) at $(0, 0)$, so this says $(\sqrt t\,I)_*\nu_I = \delta_0$, and applying $(t^{-1/2}I)_*$ gives
$\nu_I = \delta_0$. If $s > 0$, applying $(t^{-1/2}I)_*$ gives $\nu_I = (aI)_*\nu_I$ with
$a := (s/t)^{1/2} \in (0,1)$, and the fact gives $\nu_I = \delta_0$. Both contradict (K4).

\emph{(2), the gauge.} Let $\kappa_{s,t} := \pi_*\mu_{s,t}$ and $\Phi_{s,t}f := \kappa_{s,t} * f$ be
the marginal family of \ref{lem:isotropic-marginal-family}, which satisfies (A1)--(A8) and (ND) of
\prog{V}{Def.~3.1} with the same $(S_\lambda)$. \prog{V}{Prop.~6.6}, restated in the part used: \emph{for a
family of operators on $L^1(\RR)$ satisfying (A1)--(A8) and (ND) with dilation action
$(S_\lambda)$, the map $\ell(\lambda) := S_\lambda 1$ is a continuous strictly increasing bijection of
$(0,\infty)$ onto $(0,\infty)$ with $\ell(1) = 1$, and $\chi := \ell^{-1}$ on $(0,\infty)$, with
$\chi(0) := 0$, is an increasing bijection of $[0,\infty)$ satisfying
$\chi(S_\lambda t) = \lambda\,\chi(t)$.} At $t = 0$ the relation holds because an increasing
bijection of $[0,\infty)$ fixes $0$, so $S_\lambda 0 = 0$. Then $\varphi = \chi^2$ is an increasing
bijection of $[0,\infty)$, so is $\theta := \varphi^{-1}$, and
$\varphi(S_\lambda t) = \lambda^2\varphi(t)$. For $u \ge 0$ this gives
$\varphi(S_\lambda\theta(u)) = \lambda^2 u$, that is, $S_\lambda\theta(u) = \theta(\lambda^2 u)$.

\emph{(2), the kernel at $I$.} Put $m := \mu_{0,\theta(1)}$. By (I3) it is rotation invariant, and
by (ND) it is not $\delta_0$, as $\theta(1) > 0$. For $s > 0$, (I4) with $\lambda := \sqrt s$ at the
pair $(0, \theta(1))$ gives
$(\sqrt s\,I)_*m = \mu_{S_{\sqrt s}0,\,S_{\sqrt s}\theta(1)} = \mu_{0,\theta(s)}$; at $s = 0$ both
sides are $\delta_0$. So $\mu_{0,\theta(s)} = (\sqrt s\,I)_*m$ for every $s \ge 0$. The law $m$ is
self-decomposable: for $c \in (0,1)$, (I1) at $0 \le \theta(c^2) \le \theta(1)$ gives
$m = \mu_{\theta(c^2),\theta(1)} * \mu_{0,\theta(c^2)} = \mu_{\theta(c^2),\theta(1)} * (cI)_*m$, which
is \ref{eq:self-decomposable} with the residual $\rho_c := \mu_{\theta(c^2),\theta(1)}$.

\emph{(2), the family.} By \ref{prop:gw-ray-families}(3) there is exactly one
$(\GL(d), \text{rays})$-cascade family of kernels $(\nu_{\Sigma,\Sigma'})$ with $\nu_I = m$, and its
kernels from scale zero are $\nu_\Sigma = (\Sigma^{1/2})_*m$; in particular
$\nu_{uI} = (\sqrt u\,I)_*m = \mu_{0,\theta(u)}$ for $u \ge 0$. Each $\nu_\Sigma$ is
self-decomposable as an image of $m$ (\ref{prop:gw-ray-families}(1)), so no $\widehat\nu_\Sigma$ has
a zero, and by \ref{lem:gw-kernels-from-origin}(2) $\nu_{\Sigma,\Sigma'}$ is the only probability
measure $\rho$ with $\nu_{\Sigma'} = \nu_\Sigma * \rho$. For $0 \le u \le v$, (I1) at
$0 \le \theta(u) \le \theta(v)$ gives
$\nu_{vI} = \mu_{0,\theta(v)} = \mu_{\theta(u),\theta(v)} * \mu_{0,\theta(u)} = \nu_{uI} *
\mu_{\theta(u),\theta(v)}$, so $\nu_{uI,vI} = \mu_{\theta(u),\theta(v)}$. Taking $u = \varphi(s)$ and
$v = \varphi(t)$ gives $\mu_{s,t} = \nu_{\varphi(s)I,\varphi(t)I}$. If $(\nu'_{\Sigma,\Sigma'})$ is a
second $(\GL(d), \text{rays})$-cascade family with this property, then at $s = 0$, $t = \theta(1)$ it
gives $\nu'_I = \mu_{0,\theta(1)} = m$, and \ref{prop:gw-ray-families}(3) gives $\nu' = \nu$.

\emph{(2), the similarities.} First, $(\mathrm S(d), \text{rays})$ is a wedge datum, for every
$d \ge 1$. $\mathrm S(d)$ is a subgroup of $\GL(d)$ containing the dilations, as
$(cQ)(c'Q') = cc'\,QQ'$ and $(cQ)^{-1} = c^{-1}Q^{\tp}$. It is closed in $\GL(d)$: on the
similarities $AA^{\tp} = c^2I$ is a multiple of $I$, a closed condition, so an invertible limit $A$
of similarities has $AA^{\tp} = kI$ with $k \ge 0$, and $k > 0$ as $\det A \ne 0$; then
$Q := k^{-1/2}A$ is orthogonal and $A = k^{1/2}Q$. \ref{lem:gw-scheme-claims}(2) applies. Next,
$\mathrm S(d) \subseteq \GL(d)$, and the order is the same: (K1), (K2) and (K4) do not involve the
group, and (K3) for $\mathrm S(d)$ is part of (K3) for $\GL(d)$. So $(\nu_{\Sigma,\Sigma'})$ is an
$(\mathrm S(d), \text{rays})$-cascade family, with the same kernels and hence the same slice.

\emph{(3).} Put $m' := \nu'_I$, which is not $\delta_0$ by (K4). As in (1), (K3) with
$A = \sqrt u\,I$ and (K1) at $u = 0$ give $\nu'_{uI} = (\sqrt u\,I)_*m'$ for $u \ge 0$. At the pairs
$(0, t)$, the hypothesis gives $\mu_{0,t} = \nu'_{\psi(t)I} = (\sqrt{\psi(t)}\,I)_*m'$, and (2) gives
$\mu_{0,t} = \nu_{\varphi(t)I} = (\sqrt{\varphi(t)}\,I)_*m$. At $t_0 := \theta(1)$, where
$\varphi(t_0) = 1$, this reads $m = (\sqrt{\psi(t_0)}\,I)_*m'$. So for $t > 0$,
$(\sqrt{\psi(t)}\,I)_*m' = (\sqrt{\varphi(t)\psi(t_0)}\,I)_*m'$, both scalars being positive, and
applying $(\psi(t)^{-1/2}I)_*$ gives $(aI)_*m' = m'$ with
$a := (\varphi(t)\psi(t_0)/\psi(t))^{1/2}$. The fact forces $a = 1$, since $m' \ne \delta_0$. Hence
$\psi(t) = c\,\varphi(t)$ for $t > 0$ with $c := \psi(t_0) > 0$, and at $t = 0$ both sides vanish.
Conversely, let $c > 0$ and $\nu^c_{\Sigma,\Sigma'} := \nu_{c^{-1}\Sigma,\,c^{-1}\Sigma'}$. The ray
order is invariant under $\Sigma \mapsto c^{-1}\Sigma$, which is the action of $c^{-1/2}I$, so
$\nu^c$ is indexed by the same pairs. (K1) and (K2) for $\nu^c$ are those of $\nu$, the map
$(\Sigma,\Sigma') \mapsto (c^{-1}\Sigma, c^{-1}\Sigma')$ being continuous; (K3) holds because
$A \cdot (c^{-1}\Sigma) = c^{-1}(A \cdot \Sigma)$; and (K4) because
$\nu^c_{0,I} = \nu_{c^{-1}I} = (c^{-1/2}I)_*m \ne \delta_0$. Finally
$\nu^c_{c\varphi(s)I,\,c\varphi(t)I} = \nu_{\varphi(s)I,\varphi(t)I} = \mu_{s,t}$.

\emph{(4).} Two $(\GL(d), \text{rays})$-cascade families with the same slice have the same kernel
$\nu_{0,I}$, and \ref{prop:gw-ray-families}(3) makes them equal. Let $d \ge 2$ and let
$(\nu_{\Sigma,\Sigma'})$ be an $(\mathrm S(d), \text{rays})$-cascade family, $(\mathrm S(d),
\text{rays})$ being a wedge datum by (2). The kernel
$m := \nu_I$ is rotation invariant by (K3) at $Q \in \Orth(d) \subseteq \mathrm S(d)$ and the pair
$(0, I)$, self-decomposable (\ref{prop:gw-ray-families}(2)), and not $\delta_0$ by (K4); as in (1),
$\nu_{sI} = (\sqrt s\,I)_*m$ for $s \ge 0$. The two families: (i) the
$(\GL(d), \text{rays})$-cascade family with kernel $m$ at $I$ (\ref{prop:gw-ray-families}(3)), read as
an $(\mathrm S(d), \text{rays})$-cascade family as in (2), whose kernels from scale zero are
$(\Sigma^{1/2})_*m$; (ii) a family with the kernels from scale zero
$\nu^{\mathrm{ii}}_\Sigma := ((\det\Sigma)^{1/(2d)}I)_*m$. These satisfy (a)--(c) of
\ref{lem:gw-kernels-from-origin}(2) for $G = \mathrm S(d)$: (a) as $\det 0 = 0$ and $\det I = 1$;
(b) as $\det(cQ\Sigma Q^{\tp}c) = c^{2d}\det\Sigma$ and $Q_*m = m$; (c) as
$\Sigma \mapsto (\det\Sigma)^{1/(2d)}$ is continuous on $\PSD(d)$. Each is self-decomposable as an
image of $m$, so \ref{prop:gw-ray-families}(2) gives the family. Both have kernel $m$ at $I$, so on
the slice the kernels from scale zero of the three families are all $(\sqrt s\,I)_*m = \nu_{sI}$. No
transform $\widehat\nu_\Sigma$ has a zero (\ref{prop:gw-ray-families}(1)), and by
\ref{lem:gw-kernels-from-origin}(2) the increment at $(sI, tI)$ of each of the three families is the
only probability measure $\rho$ with $\nu_{tI} = \nu_{sI} * \rho$: the three slices agree. Finally, let
$\Sigma \in \PD(d)$ have $(\Sigma^{1/2})_*m = (vI)_*m$ with $v := (\det\Sigma)^{1/(2d)} > 0$. Then
$P := v^{-1}\Sigma^{1/2}$ fixes $m$, so $P \in \operatorname{Stab}(m) = \Orth(d)$ by
\ref{lem:parameter-space}(2), $m$ being rotation invariant and not $\delta_0$. As $P$ is symmetric,
$I = PP^{\tp} = v^{-2}\Sigma$, and $\Sigma = v^2I$ is a multiple of $I$.
\end{proof}

\section{The parameter space}
\label{sec:parameter-space}

The covariance of an affine family fixes its parametrization up to one constant
(\ref{lem:affine-gauge}). Why the parameter space is $\PD(d)$: a full law has a compact stabilizer,
and an orbit of $\GL(d)$ is indexed by $\PD(d)$ exactly when it has a rotation-invariant member, so
the isotropic member is an axiom (\ref{lem:parameter-space}).

% draft: section 2.2, Lemma 2.5. The draft asks for order-preserving bijections; the proof reads
% equivariance only, so the node states the stronger form.
\begin{lemma}[the affine gauge is a constant]
\label{lem:affine-gauge}
\uses{def:affine-cascade-family}
\lean{AffineHemigroup.affine_gauge, AffineHemigroup.affine_gauge_slice}\leanok
Let $h : \PD(d) \to \PD(d)$ satisfy $h(A\Sigma A^{\tp}) = A\,h(\Sigma)\,A^{\tp}$ for every
$A \in \GL(d)$ and $\Sigma \in \PD(d)$. Then there is a $c > 0$ with $h(\Sigma) = c\,\Sigma$ for
every $\Sigma$. On the isotropic slice $\{sI : s > 0\}$ the dilation $x \mapsto \lambda x$,
$\lambda > 0$, which is $A = \lambda I$, acts by $s \mapsto \lambda^2 s$.
\statusT\quad\emph{(Cites nothing. The draft's statement asks for order-preserving bijections;
neither hypothesis is used. \textbf{Not claimed}: that a family stated in some other
equivariant parametrization exists and is carried to the canonical one by such an $h$ --- that is
an abstract form of the axiom, which the draft's check item~4 leaves for a referee --- and the
identification of $s$ with the variance gauge of [V, Prop.~6.6], which is a reading: in the
Gaussian family $s$ is the variance.)}
\end{lemma}

\begin{proof}
\leanok
For $Q \in \Orth(d)$, $QQ^{\tp} = I$ gives $Q\,h(I)\,Q^{\tp} = h(I)$: the symmetric matrix
$H := h(I)$ commutes with every orthogonal matrix. Taking $Q$ the reflection
$I - 2e_ie_i^{\tp}$, the entry $(QHQ^{\tp})_{ik}$ is $-H_{ik}$ for $k \ne i$, so $H_{ik} = 0$ and
$H$ is diagonal; taking $Q$ the permutation matrix exchanging $e_i$ and $e_k$ gives
$H_{ii} = H_{kk}$. So $H = cI$, and $c > 0$ because $H \in \PD(d)$. For $\Sigma \in \PD(d)$ put
$A = \Sigma^{1/2}$: then $h(\Sigma) = h(AIA^{\tp}) = A(cI)A^{\tp} = c\,\Sigma$. Finally
$(\lambda I)(sI)(\lambda I)^{\tp} = \lambda^2 s\,I$.
\end{proof}

% not in the draft: the plan's item D1 (notes/PLAN-article-structure.md, section 4.2; Q-0265).
% Why the parameter space is PD(d): a kernel that is not carried by a proper subspace has a compact
% stabilizer, a compact group containing O(d) is O(d), and an orbit of GL(d) can be indexed by PD(d)
% exactly when it has a rotation-invariant member. What writing it down found: the brief's form, "a
% law invariant under a noncompact closed subgroup is delta_0", is false (delta_{e_1} is invariant
% under every A with A e_1 = e_1); the true statement is (1), and fullness is what nondegeneracy gives
% only when G contains O(d) (clause (5)).
% CHANGED Q-0281: typed (the statement skeleton). Clause (3) is split to lem:parameter-space-indexing
% and clauses (4)-(5) to lem:parameter-space-examples (ADR-0025; records/formalization/PRICING.md:
% stabilizers, indexing and examples are three arguments of different kind and cost); the text is
% unchanged and the \lean tag names the children's declarations too.
\begin{lemma}[the parameter space: compact stabilizers and the isotropic member]
\label{lem:parameter-space}
\uses{def:gw-cascade-family,lem:affine-gauge}
\lean{AffineHemigroup.stab_isClosedSubgroupGL, AffineHemigroup.stab_isCompact_of_isFull,
AffineHemigroup.carried_by_subspace_of_not_isCompact, AffineHemigroup.stab_dirac_ebasis,
AffineHemigroup.isFull_of_isRotationInvariant, AffineHemigroup.parameter_space_tfae,
AffineHemigroup.parameter_space_index_bijOn, AffineHemigroup.parameter_space_index_unique,
AffineHemigroup.parameter_space_atoms, AffineHemigroup.parameter_space_scheme,
AffineHemigroup.gram_glSet_eq_pdSet, AffineHemigroup.parameter_space_diag_example}\leanok
Let $d \ge 1$. For a probability measure $\mu$ on $\RR^d$ write
$\operatorname{Stab}(\mu) := \{A \in \GL(d) : A_*\mu = \mu\}$ and
$\operatorname{Orb}(\mu) := \{A_*\mu : A \in \GL(d)\}$, and call $\mu$ \emph{full} if
$\mu(V) < 1$ for every linear subspace $V \ne \RR^d$.
\begin{enumerate}
\item \emph{(Compact stabilizers.)} $\operatorname{Stab}(\mu)$ is a closed subgroup of $\GL(d)$. If
$\mu$ is full, $\operatorname{Stab}(\mu)$ is compact. So a law whose stabilizer is not compact is
carried by a proper linear subspace; for $d \ge 2$ it need not be $\delta_0$: the stabilizer of
$\delta_{e_1}$ contains every $A$ with $Ae_1 = e_1$.
\item \emph{(Rotation-invariant laws.)} If $Q_*\mu = \mu$ for every $Q \in \Orth(d)$ and
$\mu \ne \delta_0$, then $\mu$ is full and $\operatorname{Stab}(\mu) = \Orth(d)$.
\item \emph{(The isotropic member is the axiom of the parameter space $\PD(d)$.)} Let $\mu_0$ be
full. The following are equivalent:
(a) $\operatorname{Orb}(\mu_0)$ has a rotation-invariant member;
(b) $\operatorname{Stab}(\mu_0) = C\,\Orth(d)\,C^{-1}$ for some $C \in \GL(d)$;
(c) there is a map $\Sigma \mapsto \mu_\Sigma$ from $\PD(d)$ into $\operatorname{Orb}(\mu_0)$ with
$A_*\mu_\Sigma = \mu_{A\Sigma A^{\tp}}$ for every $A \in \GL(d)$ and $\Sigma \in \PD(d)$.
When they hold, every map as in (c) is a bijection of $\PD(d)$ onto $\operatorname{Orb}(\mu_0)$, and
for any two such maps $\mu$ and $\mu'$ there is a $c > 0$ with $\mu'_\Sigma = \mu_{c\Sigma}$ for
every $\Sigma \in \PD(d)$.
\item \emph{(An orbit without an isotropic member.)} Let $p_1, \dots, p_d > 0$ be distinct with
$p_1 + \dots + p_d = 1$, and $\mu_0 := \sum_jp_j\,\delta_{e_j}$. Then $\mu_0$ is full,
$\operatorname{Stab}(\mu_0) = \{I\}$, the map $A \mapsto A_*\mu_0$ is a bijection of $\GL(d)$ onto
$\operatorname{Orb}(\mu_0)$, and no member of $\operatorname{Orb}(\mu_0)$ is rotation invariant.
\item \emph{(In the scheme.)} Let $(G, W)$ be a wedge datum with $\Orth(d) \subseteq G$ and
$(\mu_{\Sigma,\Sigma'})$ a $(G, W)$-cascade family of kernels (\ref{def:gw-cascade-family}(2)). Then
$\mu_I$ is rotation invariant and full, $\operatorname{Stab}(\mu_I) = \Orth(d)$, and
$\Sigma \mapsto \mu_\Sigma$ is injective on $\{AA^{\tp} : A \in G\}$, which is $\PD(d)$ when
$G = \GL(d)$. Without $\Orth(d) \subseteq G$, (K4) need not make $\mu_I$ full: for $d \ge 2$, $G$
the invertible diagonal matrices and $W$ the Loewner order, the centred Gaussian laws
$\mu_{\Sigma,\Sigma'}$ with $\widehat\mu_{\Sigma,\Sigma'}(\omega) =
\exp[-\tfrac12(\Sigma' - \Sigma)_{11}\,\omega_1^2]$ form a $(G, W)$-cascade family of kernels whose
$\mu_I$ is carried by the first axis and has a stabilizer that is not compact.
\end{enumerate}
\statusT\quad\emph{(Cites nothing beyond the uniqueness of the Fourier transform of a probability
measure on $\RR^d$, the polar decomposition of an invertible matrix, and dominated convergence.
\textbf{What it is worth.} It says why the classical theory's parameter space is $\PD(d)$ and what
that costs. A nondegenerate kernel of a family covariant under a group containing $\Orth(d)$ is
full, so its stabilizer is compact, and being rotation invariant it is exactly $\Orth(d)$; the
kernels from scale zero then correspond one to one to $\GL(d)/\Orth(d) = \PD(d)$. For an orbit of
$\GL(d)$ given as a set of laws, indexing by $\PD(d)$ is possible exactly when the orbit has an
isotropic member, and then unique up to the constant of \ref{lem:affine-gauge}: clause (3)'s last
sentence is that lemma, and the two fit as existence and uniqueness of the index. In
\ref{def:gw-cascade-family} the isotropic member is built into the indexing (its annotation); clause
(4) shows that it is an axiom and not a consequence, since a full law with trivial stabilizer has an
orbit in bijection with $\GL(d)$ itself, a parameter space of dimension $d^2$ against
$d(d+1)/2$ for $\PD(d)$. \textbf{Against the brief.} The plan's form ``a law invariant under a
noncompact closed subgroup of $\GL(d)$ is $\delta_0$'' is false as stated ($\delta_{e_1}$ in (1)); the
correct statement is (1)'s, and it gives $\delta_0$ when the subgroup also fixes no proper subspace
carrying the law, which for a rotation-invariant law is (2). \textbf{Not claimed}: that an orbit
$\GL(d)/K$ is a manifold of dimension $d^2 - \dim K$ for a general compact $K$ (the closed-subgroup
theorem, not used); a classification of the possible stabilizers between $\{I\}$ and $\Orth(d)$; and
anything about stabilizers in groups $G$ not containing $\Orth(d)$ beyond the example of (5).)}
\end{lemma}

\begin{proof}
\leanok
\emph{(1).} By uniqueness of the transform, $A_*\mu = \mu$ exactly when
$\widehat\mu(A^{\tp}\omega) = \widehat\mu(\omega)$ for every $\omega$. As $(AB)_*\mu = A_*(B_*\mu)$,
$\operatorname{Stab}(\mu)$ is closed under products, and $A_*\mu = \mu$ gives
$\mu = (A^{-1})_*\mu$; so it is a subgroup. If $A_n \to A$ in $\GL(d)$ with
$A_n \in \operatorname{Stab}(\mu)$, then $\widehat\mu(A^{\tp}\omega) = \lim_n\widehat\mu(A_n^{\tp}\omega)
= \widehat\mu(\omega)$ by continuity of $\widehat\mu$; so the subgroup is closed in $\GL(d)$.

Let $\mu$ be full. The function $f(B) := \int\min(1, |Bx|)\,\mu(dx)$ on the linear maps of $\RR^d$
is continuous, by dominated convergence, the integrand being at most $1$. On the unit sphere
$\{\|B\| = 1\}$, which is compact, it attains its minimum $c$ at some $B_0$, and $c > 0$: if
$f(B_0) = 0$, then $B_0x = 0$ for $\mu$-almost every $x$, and $\mu$ is carried by $\ker B_0$, a
proper subspace as $B_0 \ne 0$, against fullness. Also
$g(t) := \int\min(1, |y|/t)\,\mu(dy) \to 0$ as $t \to \infty$, by dominated convergence again; take
$T \ge 1$ with $g(t) < c$ for $t \ge T$. If $A \in \operatorname{Stab}(\mu)$ had $t := \|A\| > T$, then
$B := A/t$ lies on the unit sphere and, since $A_*\mu = \mu$,
\[
  c \le f(B) = \int\min\bigl(1, |Ax|/t\bigr)\,\mu(dx) = \int\min\bigl(1, |y|/t\bigr)\,\mu(dy) = g(t) < c,
\]
which is absurd. So $\|A\| \le T$ on $\operatorname{Stab}(\mu)$, which therefore lies in a compact
set $K$ of matrices. On $K$ the continuous function $|\det|$ is bounded by some $C \ge 1$; applied
to $A^{-1}$, which is in the group, this gives $|\det A| \ge 1/C$ on $\operatorname{Stab}(\mu)$, and
so on its closure among matrices. A matrix of that closure is therefore invertible, and lies in
$\operatorname{Stab}(\mu)$ because the subgroup is closed in $\GL(d)$. So $\operatorname{Stab}(\mu)$
is a closed subset of the compact set $K$, hence compact. For $\delta_{e_1}$,
$A_*\delta_{e_1} = \delta_{Ae_1}$, and for $d \ge 2$ the matrices $I + t\,e_1e_2^{\tp}$, $t \in \RR$,
fix $e_1$ and are unbounded.

\emph{(2).} Suppose $\mu(V) = 1$ for a proper subspace $V$, and choose a unit vector $u \perp V$, so
that the hyperplane $u^\perp$ carries $\mu$. For $Q \in \Orth(d)$,
$\mu(Q(u^\perp)) = (Q^{-1})_*\mu(u^\perp) = \mu(u^\perp) = 1$, and $Q(u^\perp) = (Qu)^\perp$; every unit
vector is some $Qu$, so each coordinate hyperplane $e_j^\perp$ carries $\mu$, and so does their
intersection, a finite intersection of sets of full measure. That intersection is $\{0\}$, so
$\mu = \delta_0$, a contradiction: $\mu$ is full. By (1), $K := \operatorname{Stab}(\mu)$ is compact,
and it contains $\Orth(d)$. Let $A \in K$, and write $A = PQ$ with $P := (AA^{\tp})^{1/2} \in \PD(d)$
and $Q := P^{-1}A$, which is orthogonal since $QQ^{\tp} = P^{-1}AA^{\tp}P^{-1} = I$. Then
$P = AQ^{-1} \in K$, so all powers $P^n$, $n \in \ZZ$, lie in the compact set $K$; an eigenvalue
$p > 0$ of $P$ has $p^n$ bounded over $n \in \ZZ$, so $p = 1$, $P = I$ and $A = Q \in \Orth(d)$. So
$K = \Orth(d)$.

\emph{(3).} Two remarks first. If $\mu$ is full and $B \in \GL(d)$, then $B_*\mu$ is full, since
$(B_*\mu)(V) = \mu(B^{-1}V)$; and $\operatorname{Stab}(B_*\mu) = B\,\operatorname{Stab}(\mu)\,B^{-1}$,
since $A_*B_*\mu = B_*\mu$ exactly when $(B^{-1}AB)_*\mu = \mu$.
(a)$\Rightarrow$(b): let $\nu := B_*\mu_0$ be rotation invariant. It is full, hence not $\delta_0$, and
by (2) $\operatorname{Stab}(\nu) = \Orth(d)$; so $\operatorname{Stab}(\mu_0) = B^{-1}\Orth(d)B$.
(b)$\Rightarrow$(c): put $\nu := (C^{-1})_*\mu_0$, whose stabilizer is $C^{-1}\operatorname{Stab}(\mu_0)C
= \Orth(d)$, and $\mu_\Sigma := (\Sigma^{1/2})_*\nu \in \operatorname{Orb}(\mu_0)$. For $A \in \GL(d)$
and $\Sigma \in \PD(d)$ let $\Sigma' := A\Sigma A^{\tp}$ and $Q := \Sigma'^{-1/2}A\Sigma^{1/2}$; then
$QQ^{\tp} = \Sigma'^{-1/2}A\Sigma A^{\tp}\Sigma'^{-1/2} = I$, so $Q$ is orthogonal, $A\Sigma^{1/2} =
\Sigma'^{1/2}Q$, and $A_*\mu_\Sigma = (\Sigma'^{1/2})_*Q_*\nu = \mu_{\Sigma'}$.
(c)$\Rightarrow$(a): for $Q \in \Orth(d)$, $Q_*\mu_I = \mu_{QQ^{\tp}} = \mu_I$, and $\mu_I$ is a
member of the orbit.

Now let $\mu$ be a map as in (c). Its member $\mu_I$ is rotation invariant and, lying in the orbit of
a full law, full; by (2) $\operatorname{Stab}(\mu_I) = \Orth(d)$. Covariance with $A = \Sigma^{1/2}$
gives $\mu_\Sigma = (\Sigma^{1/2})_*\mu_I$. \emph{Onto}: $\operatorname{Orb}(\mu_0) =
\operatorname{Orb}(\mu_I)$, and $A_*\mu_I = \mu_{AA^{\tp}}$. \emph{One to one}: if $\mu_\Sigma =
\mu_{\Sigma'}$, then $(\Sigma'^{-1/2}\Sigma^{1/2})_*\mu_I = \mu_I$, so $Q := \Sigma'^{-1/2}\Sigma^{1/2}$ is
orthogonal and $\Sigma = \Sigma^{1/2}(\Sigma^{1/2})^{\tp} = \Sigma'^{1/2}QQ^{\tp}\Sigma'^{1/2} = \Sigma'$. For a
second map $\mu'$, $h := \mu^{-1} \circ \mu'$ is a map of $\PD(d)$ into itself, and
$\mu_{h(A\Sigma A^{\tp})} = \mu'_{A\Sigma A^{\tp}} = A_*\mu'_\Sigma = A_*\mu_{h(\Sigma)} = \mu_{Ah(\Sigma)A^{\tp}}$;
injectivity gives $h(A\Sigma A^{\tp}) = A\,h(\Sigma)\,A^{\tp}$, and \ref{lem:affine-gauge} gives
$h(\Sigma) = c\,\Sigma$ with $c > 0$.

\emph{(4).} A proper subspace misses some $e_j$, since otherwise it would contain their span, and then
it has $\mu_0$-measure at most $1 - p_j < 1$: $\mu_0$ is full. For $A \in \GL(d)$,
$A_*\mu_0 = \sum_jp_j\,\delta_{Ae_j}$ with the points $Ae_j$ distinct. If $A_*\mu_0 = \mu_0$, the atom of
mass $p_j$, which is unique since the $p_j$ are distinct, sits at $e_j$ and at $Ae_j$, so $Ae_j = e_j$
for every $j$ and $A = I$. Hence $A_*\mu_0 = B_*\mu_0$ only when $B^{-1}A = I$, which is the bijection.
A member $A_*\mu_0$ gives mass $p_1 > 0$ to the nonzero point $x := Ae_1$, and gives positive mass
only to the $d$ points $Ae_j$. Suppose it were rotation invariant, and let $r := |x|$. Every $y$
with $|y| = r$ is $Qx$ for an orthogonal $Q$ (the reflection in the hyperplane orthogonal to
$x - y$ when $y \ne x$), so it carries the same positive mass as $x$ and is one of the points
$Ae_j$. The $2d$ distinct points $\pm r\,e_j$ all have norm $r$, and would thus be among $d$
points, which is impossible.

\emph{(5).} For $Q \in \Orth(d) \subseteq G$, $Q \cdot 0 = 0$ and $Q \cdot I = I$, so (K3) gives
$Q_*\mu_I = \mu_I$; by (K4) $\mu_I \ne \delta_0$, and (2) makes it full with stabilizer $\Orth(d)$.
For $A \in G$, (K3) gives $\mu_{AA^{\tp}} = \mu_{A\cdot 0,\,A\cdot I} = A_*\mu_I$, and
$A_*\mu_I = B_*\mu_I$ forces $B^{-1}A \in \Orth(d)$, hence $AA^{\tp} = BB^{\tp}$. Every
$\Sigma \in \PD(d)$ is $AA^{\tp}$ with $A = \Sigma^{1/2} \in \GL(d)$. For the example: the
invertible diagonal matrices form a closed subgroup containing the dilations, and the Loewner order is
invariant under every $A \in \GL(d)$, as $A(\Sigma' - \Sigma)A^{\tp}$ is positive semidefinite with
$\Sigma' - \Sigma$, and contains the rays, as $c'\Sigma - c\Sigma = (c' - c)\Sigma$; so
$(G, W)$ is a wedge datum. For $\Sigma \le \Sigma'$, $(\Sigma' - \Sigma)_{11} \ge 0$, so the
transform is that of the image of the centred normal law of variance $(\Sigma' - \Sigma)_{11}$ under
$t \mapsto te_1$. The exponents are additive in $\Sigma' - \Sigma$ and vanish at $\Sigma = \Sigma'$,
which with uniqueness of the transform is (K1); they are continuous, which is (K2); for
$A = \diag(a_1,\dots,a_d)$, $\widehat\mu_{\Sigma,\Sigma'}(A\omega) = \exp[-\tfrac12a_1^2(\Sigma' -
\Sigma)_{11}\omega_1^2]$ and $(A\Sigma'A - A\Sigma A)_{11} = a_1^2(\Sigma' - \Sigma)_{11}$, which is (K3);
and $\widehat\mu_{0,I}(e_1) = e^{-1/2} \ne 1$, which is (K4). Its $\mu_I$ is carried by the first
axis, and its stabilizer contains the unbounded matrices $\diag(1, t, 1, \dots, 1)$, $t > 0$.
\end{proof}

% CHANGED Q-0281: split from lem:parameter-space (ADR-0025; records/formalization/PRICING.md): clause
% (3), the indexing of an orbit by PD(d), copied verbatim with the notation it reads. The parent's text
% and numbering are unchanged.
\begin{lemma}[the parameter space: indexing an orbit by $\PD(d)$]
\label{lem:parameter-space-indexing}
\uses{lem:parameter-space,lem:affine-gauge}
\lean{AffineHemigroup.parameter_space_tfae, AffineHemigroup.parameter_space_index_bijOn,
AffineHemigroup.parameter_space_index_unique}\leanok
Let $d \ge 1$, and use the notation of \ref{lem:parameter-space}. Let $\mu_0$ be full. The
following are equivalent:
(a) $\operatorname{Orb}(\mu_0)$ has a rotation-invariant member;
(b) $\operatorname{Stab}(\mu_0) = C\,\Orth(d)\,C^{-1}$ for some $C \in \GL(d)$;
(c) there is a map $\Sigma \mapsto \mu_\Sigma$ from $\PD(d)$ into $\operatorname{Orb}(\mu_0)$ with
$A_*\mu_\Sigma = \mu_{A\Sigma A^{\tp}}$ for every $A \in \GL(d)$ and $\Sigma \in \PD(d)$.
When they hold, every map as in (c) is a bijection of $\PD(d)$ onto $\operatorname{Orb}(\mu_0)$, and
for any two such maps $\mu$ and $\mu'$ there is a $c > 0$ with $\mu'_\Sigma = \mu_{c\Sigma}$ for
every $\Sigma \in \PD(d)$.
\statusT\quad\emph{(Cites nothing. Clause (3) of \ref{lem:parameter-space}, split off on 2026-10-04
(Q-0281, ADR-0025); the numbering stays the parent's.)}
\end{lemma}

\begin{proof}
\leanok
This is the part of the proof of \ref{lem:parameter-space} headed (3), which reads that node's
clause (2) and \ref{lem:affine-gauge}.
\end{proof}

% CHANGED Q-0281: split from lem:parameter-space (ADR-0025; records/formalization/PRICING.md): clauses
% (4) and (5), the two examples and the reading in the scheme, copied verbatim. The parent's text and
% numbering are unchanged.
\begin{lemma}[the parameter space: an orbit without an isotropic member, and the scheme]
\label{lem:parameter-space-examples}
\uses{lem:parameter-space,def:gw-cascade-family,lem:gw-scheme-claims}
\lean{AffineHemigroup.parameter_space_atoms, AffineHemigroup.parameter_space_scheme,
AffineHemigroup.gram_glSet_eq_pdSet, AffineHemigroup.parameter_space_diag_example}\leanok
Let $d \ge 1$, and use the notation of \ref{lem:parameter-space}.
\begin{enumerate}
\item[(4)] \emph{(An orbit without an isotropic member.)} Let $p_1, \dots, p_d > 0$ be distinct with
$p_1 + \dots + p_d = 1$, and $\mu_0 := \sum_jp_j\,\delta_{e_j}$. Then $\mu_0$ is full,
$\operatorname{Stab}(\mu_0) = \{I\}$, the map $A \mapsto A_*\mu_0$ is a bijection of $\GL(d)$ onto
$\operatorname{Orb}(\mu_0)$, and no member of $\operatorname{Orb}(\mu_0)$ is rotation invariant.
\item[(5)] \emph{(In the scheme.)} Let $(G, W)$ be a wedge datum with $\Orth(d) \subseteq G$ and
$(\mu_{\Sigma,\Sigma'})$ a $(G, W)$-cascade family of kernels (\ref{def:gw-cascade-family}(2)). Then
$\mu_I$ is rotation invariant and full, $\operatorname{Stab}(\mu_I) = \Orth(d)$, and
$\Sigma \mapsto \mu_\Sigma$ is injective on $\{AA^{\tp} : A \in G\}$, which is $\PD(d)$ when
$G = \GL(d)$. Without $\Orth(d) \subseteq G$, (K4) need not make $\mu_I$ full: for $d \ge 2$, $G$
the invertible diagonal matrices and $W$ the Loewner order, the centred Gaussian laws
$\mu_{\Sigma,\Sigma'}$ with $\widehat\mu_{\Sigma,\Sigma'}(\omega) =
\exp[-\tfrac12(\Sigma' - \Sigma)_{11}\,\omega_1^2]$ form a $(G, W)$-cascade family of kernels whose
$\mu_I$ is carried by the first axis and has a stabilizer that is not compact.
\end{enumerate}
\statusT\quad\emph{(Cites nothing beyond the uniqueness of the Fourier transform. Clauses (4) and
(5) of \ref{lem:parameter-space}, split off on 2026-10-04 (Q-0281, ADR-0025); the numbering stays
the parent's. The Lean types (5)'s example with the laws written as Gaussian measures of covariance
$\diag((\Sigma' - \Sigma)_{11}, 0, \dots, 0)$ and states the displayed transform as a conjunct.)}
\end{lemma}

\begin{proof}
\leanok
This is the part of the proof of \ref{lem:parameter-space} headed (4) and (5), which reads that
node's clauses (1) and (2); that the Loewner order is invariant and contains the rays is
\ref{lem:gw-scheme-claims}(2).
\end{proof}

\section{Infinitely divisible increments from the cascade}
\label{sec:null-array}

Along an order with the path condition, every increment of a cascade is infinitely divisible, by
Khintchine's null-array theorem (\ref{thm:cascade-infinitely-divisible}). The rigidity of
\S\ref{sec:affine-single-step} does not use it.

% not in the draft as a node: the draft's Lemma 2.6, last sentence, by the null-array argument of
% [V, Thm. 5.2] in R^d; the plan's item D4 (notes/PLAN-article-structure.md, "What the null array
% changes", the author's decision 2 of 2026-10-03; Q-0265). Ledger A6 (Sato Thm. 9.3, p. 47, read at
% the page, records/affine/READING-sato1999levy-null-array.md). The infinitesimality of the array is
% proved here from the continuity of the transforms; Sato's E 12.12 (an exercise) and Lemma 9.6 /
% Thm. 9.7 (additive processes) are read and not cited.
% CHANGED Q-0281: typed (the statement skeleton). Clause (3), the chains, is split to
% lem:path-condition-orders (ADR-0025; records/formalization/PRICING.md: matrix analysis to Lean core,
% read by thm:cone-orders, where (1), (2) and (4) spend A6); the text is unchanged and the \lean tag
% names the child's declarations too.
% CHANGED Q-0281: \uses re-pointed to the split child it reads (prop:stable-step-cone(2), the order, only).
% CHANGED Q-0281: blind review R1-2: the edge to the split child runs parent -> child (the [T] child
% no longer \uses this [A] node), as for the other [T] children of [A] parents; and R1-6: the tag
% names the declarations of lem:gw-scheme-claims(2) that make the Loewner and the ray order on PSD(d)
% partial orders, which clause (3)(a)-(b)'s "satisfy the hypothesis of (2)" includes.
% CHANGED Q-0281 wave 2: the \statusA annotation inside this statement says how the infinitesimality
% of the array is now proved (Levy's continuity theorem on R^d and the portmanteau inequality, both in
% Mathlib, no ledger entry), with the proof below; the author's decision of 2026-10-05 accepts that
% route. The statement text is unchanged. Marked here on blind review W2C-4.
% CHANGED Q-0327: notation only (whole-paper review flag 16): the $W$-chain is $\zeta$, no longer
% $\gamma$ (a stable index, the Matern index, a causal curve), here and in lem:path-condition-orders.
% CHANGED round-2 review (Gemini), item 12, the author's decision D7 (2026-11-03): the exponent cone K_alpha
% is written frak k_alpha (a subset of gl(d), in fraktur as gl(d) is), apart from the step cone
% cal K_alpha. Notation only; no claim, \uses or \lean changes.
\begin{theorem}[a continuous cascade has infinitely divisible increments]
\label{thm:cascade-infinitely-divisible}
\uses{def:gw-cascade-family,def:isotropic-kernel-family,lem:stable-step-order,lem:path-condition-orders,thm:cited-A6}
\lean{AffineHemigroup.cascade_interval_infinitelyDivisible,
AffineHemigroup.chain_infinitelyDivisible, AffineHemigroup.gw_family_infinitelyDivisible,
AffineHemigroup.isotropic_family_infinitelyDivisible, AffineHemigroup.chain_loewner,
AffineHemigroup.chain_ray_zero, AffineHemigroup.chain_ray, AffineHemigroup.chain_stepSet,
AffineHemigroup.stepSet_loewner, AffineHemigroup.stepSet_ray,
AffineHemigroup.stepSet_stable, AffineHemigroup.isWedgeDatum_loewner,
AffineHemigroup.isWedgeDatum_ray, AffineHemigroup.glSet_isClosedSubgroupGL}\leanok
Let $d \ge 1$.
\begin{enumerate}
\item \emph{(A cascade along an interval.)} Let $(\mu_{s,t})_{0 \le s \le t \le 1}$ be probability
measures on $\RR^d$ with $\mu_{s,s} = \delta_0$, $\mu_{r,t} = \mu_{s,t} * \mu_{r,s}$ whenever
$r \le s \le t$, and $(s,t) \mapsto \widehat\mu_{s,t}(\omega)$ continuous on
$\{0 \le s \le t \le 1\}$ for every $\omega \in \RR^d$. Then every $\mu_{s,t}$ is infinitely
divisible.
\item \emph{(The path condition.)} Let $\mathcal P \subseteq \PSD(d)$, let $\le_W$ be a partial order
on $\mathcal P$, and let $(\mu_{\Sigma,\Sigma'})$ be probability measures on $\RR^d$, indexed by the
pairs $\Sigma \le_W \Sigma'$ in $\mathcal P$, satisfying (K1) and (K2) of
\ref{def:gw-cascade-family} on these pairs. Call a continuous $\zeta : [0,1] \to \mathcal P$ with
$\zeta(s) \le_W \zeta(t)$ for all $s \le t$ a \emph{$W$-chain} from $\zeta(0)$ to $\zeta(1)$. If
$\Sigma$ and $\Sigma'$ are joined by a $W$-chain, then $\mu_{\Sigma,\Sigma'}$ is infinitely divisible.
In particular, for a wedge datum $(G, W)$ satisfying
\begin{enumerate}
\item[(P)] every pair $\Sigma \le_W \Sigma'$ in $\PSD(d)$ is joined by a $W$-chain,
\end{enumerate}
every $\mu_{\Sigma,\Sigma'}$ of a $(G, W)$-cascade family of kernels is infinitely divisible, and so
is every kernel from scale zero $\mu_\Sigma$.
\item \emph{(Orders with the path condition.)} The following satisfy the hypothesis of (2) for every
pair, the chain being written out.
(a) The Loewner order on $\PSD(d)$: $\zeta(s) := \Sigma + s(\Sigma' - \Sigma)$.
(b) The ray order on $\PSD(d)$ (\ref{def:gw-cascade-family}): $\zeta(s) := s\,\Sigma'$ from
$\Sigma = 0$, and $\zeta(s) := (1 + s(c - 1))\,\Sigma$ from $\Sigma$ to $\Sigma' = c\,\Sigma$.
(c) Let $\mathcal K \subseteq \PD(d)$ satisfy $QMQ^{\tp} \in \mathcal K$ and $M^\tau \in \mathcal K$
whenever $M \in \mathcal K$, $Q \in \Orth(d)$ and $\tau \in [0,1]$. On $\mathcal P = \PD(d)$ write
$\Sigma \le_{\mathcal K} \Sigma'$ when $\Sigma^{-1/2}\Sigma'\Sigma^{-1/2} \in \mathcal K$, and suppose
that this relation is a partial order. Then it satisfies the hypothesis of (2) for every pair, with
$\zeta(s) := \Sigma^{1/2}M^s\Sigma^{1/2}$ and $M := \Sigma^{-1/2}\Sigma'\Sigma^{-1/2}$.
This covers the Loewner order on $\PD(d)$ ($\mathcal K = \{M : M \ge I\}$), the ray order on $\PD(d)$
($\mathcal K = \{cI : c \ge 1\}$), and for $0 < \alpha < 2$ the order of the stable step cone,
$\mathcal K_\alpha := \{M \in \PD(d) : \lambda_{\min}(M)^{d+\alpha} \ge \det M\}$ with $\lambda_{\min}(M)$ the
least eigenvalue; $\mathcal K_\alpha$ is the set of $e^X$ with $X$ symmetric and
$(d+\alpha)\lambda_{\min}(X) \ge \tr X$. These three relations are partial orders.
\item \emph{(Isotropic families.)} Every $\mu_{s,t}$ of an isotropic cascade family in dimension $d$
(\ref{def:isotropic-kernel-family}) is infinitely divisible.
\end{enumerate}
\statusA\quad\ledger{A6}\quad\emph{(\textbf{Assignment.} \ledger{A6}, Khintchine's theorem
(\sscite{sato1999levy} Thm.~9.3, p.~47), is spent once, in (1), on an array whose row products are
all equal to the law in question, so that its convergence hypothesis is trivial. The infinitesimality
of the array is proved here from the continuity of the transforms, through \Levy's continuity theorem
on $\RR^d$ and the portmanteau inequality, both machine-checked in Mathlib (no ledger entry); (2)--(4)
are reductions to (1) and cite nothing further. \textbf{What it is worth.} It is the draft's
Lemma~2.6, last sentence, which the module had avoided until the author admitted the null array
(\texttt{notes/PLAN-article-structure.md}, decision 2): along a path of comparable parameters the
cascade cuts an increment into $n$ steps, and continuity on the compact path makes them a null array.
Neither $G$-covariance nor nondegeneracy is used, only (K1), (K2) and the chain; so every corner of
the scheme whose order has the path condition has infinitely divisible increments, the classical
corner $(\GL(d), \text{Loewner})$, the similarities with the Loewner order and the ray corners among
them. For the stable step cone the chain is the geodesic $\Sigma^{1/2}M^s\Sigma^{1/2}$, whose
infinitesimal steps $\zeta^{-1/2}\zeta'\zeta^{-1/2}$ are orthogonal conjugates of $\log M$, so they lie
in $\mathfrak k_\alpha$; that $\preceq_\alpha$ is the whole order generated by the cone field $\mathfrak k_\alpha$ (the
plan's item D3) is \ref{prop:stable-order-generated}(1), and the closed orders of (3)(c) are exactly the
orders of invariant cone fields (\ref{thm:cone-orders}(2)). Clause (4) gives the isotropic families' infinite divisibility directly;
the isotropic nodes reach it only in the reading (2.1), through \ref{thm:isotropic-classification}
and entry A4. The set $\mathcal K_\alpha$ of (3)(c) is read off the symmetric matrices of the cone
$\mathfrak k_\alpha$ of \ref{prop:stable-exponent-cone}: its defining inequality \ref{eq:stable-exponent-cone}
at a symmetric $B = X$; the clause does not use that node, and takes the fact that
$\preceq_\alpha$ is a partial order from \ref{prop:stable-step-cone}(2). At $(\GL(d), \text{Loewner})$ and
$d \ge 2$ the conclusion is weaker than \ref{thm:affine-rigidity}, which makes every increment
Gaussian without the null array. \textbf{Not claimed}:
anything for an order without the path condition, where a pair need not be the end of a continuous
cascade; any statement about a single step $\mu_{\Sigma'} = \mu_\Sigma * \rho$ that is not embedded in
a cascade (the Gaussian--\Matern{} step of \texttt{records/affine/check-gaussian-matern-step.py} is
such a step, positive and not infinitely divisible); the triplet of an increment, which is entry
A3's.)}
\end{theorem}

% CHANGED Q-0281 wave 2: proved (AffineHemigroup.CascadeDivisible with AffineHemigroup.PathCondition;
% (1), (2), (4) print A6 and Lean core, guarded). Proof of (1) rewritten to the Lean route: the
% rows are infinitesimal by a uniform small-increment statement, proved by contradiction through
% Levy's continuity theorem and the portmanteau inequality (both in Mathlib), in place of the
% coordinatewise truncation inequality; the partition has n + 1 steps. (3) points at the split child.
\begin{proof}
\leanok
\emph{Small increments.} Under the hypotheses of (1), for every $\varepsilon > 0$ and $\eta > 0$ there is
$\delta > 0$ with $\mu_{a,b}(\{|x| > \varepsilon\}) < \eta$ whenever $0 \le a \le b \le 1$ and
$b - a < \delta$. Suppose not: then there are pairs $0 \le a_n \le b_n \le 1$ with
$b_n - a_n < 1/(n+1)$ and $\mu_{a_n,b_n}(\{|x| > \varepsilon\}) \ge \eta$. For each $\omega$ the
function $f(a,b) := \widehat\mu_{a,b}(\omega)$ is continuous on the compact triangle
$\{0 \le a \le b \le 1\}$, hence uniformly continuous, and $f(a,a) = \widehat\delta_0(\omega) = 1$;
as the distance of $(a_n, b_n)$ to $(a_n, a_n)$ is $b_n - a_n \to 0$, $f(a_n, b_n) \to 1$. So the
transforms of $\mu_{a_n,b_n}$ converge pointwise to that of $\delta_0$, and by \Levy's continuity
theorem on $\RR^d$ the laws converge weakly to $\delta_0$. The portmanteau inequality for the closed
set $F := \{|x| \ge \varepsilon\}$ bounds the upper limit of $\mu_{a_n,b_n}(F)$ by $\delta_0(F) = 0$, so
$\mu_{a_n,b_n}(\{|x| > \varepsilon\}) \le \mu_{a_n,b_n}(F) < \eta$ for some $n$, a contradiction.

\emph{(1).} Fix $s \le t$, and for $n \ge 0$ put $t_{n,k} := s + k(t - s)/(n+1)$ and
$\mu_{n,k} := \mu_{t_{n,k},\,t_{n,k+1}}$, $0 \le k \le n$. For any points
$\tau_0 \le \tau_1 \le \dots \le \tau_N$ of $[0,1]$ the cascade identity and the commutativity of
convolution give, by induction on $N$,
$\mu_{\tau_0,\tau_1} * \mu_{\tau_1,\tau_2} * \dots * \mu_{\tau_{N-1},\tau_N} = \mu_{\tau_0,\tau_N}$
(for $N = 0$ the empty convolution $\delta_0 = \mu_{\tau_0,\tau_0}$; the inductive step is
$\mu_{\tau_0,\tau_1} * \mu_{\tau_1,\tau_N} = \mu_{\tau_1,\tau_N} * \mu_{\tau_0,\tau_1} = \mu_{\tau_0,\tau_N}$). So
\[
  \mu_{n,0} * \mu_{n,1} * \dots * \mu_{n,n} = \mu_{s,t} \qquad \text{for every } n .
\]
\emph{The rows are infinitesimal.} Given $\varepsilon > 0$ and $\eta > 0$, take $\delta$ from the
small increments; once $(t - s)/(n+1) < \delta$, every $\mu_{n,k}(\{|x| > \varepsilon\})$ is below
$\eta$, so $\max_k \mu_{n,k}(\{|x| > \varepsilon\}) \to 0$. \emph{Khintchine.} The
measures $\mu_{n,k}$, $0 \le k \le n$, form a null array in the measure form of \ledger{A6}, the row
products are all $\mu_{s,t}$, and without centring (the vectors $b_n$ of the entry all $0$) they
converge to $\mu_{s,t}$; so $\mu_{s,t}$ is infinitely divisible.

\emph{(2).} Let $\zeta$ be a $W$-chain from $\Sigma$ to $\Sigma'$ and put
$\nu_{u,v} := \mu_{\zeta(u),\zeta(v)}$ for $0 \le u \le v \le 1$, defined since $\zeta(u) \le_W
\zeta(v)$. (K1) at the pairs of the chain gives $\nu_{u,u} = \delta_0$ and $\nu_{u,w} = \nu_{v,w} *
\nu_{u,v}$ for $u \le v \le w$. The map $(u,v) \mapsto (\zeta(u), \zeta(v))$ is continuous into the
set of pairs, so (K2) makes $(u,v) \mapsto \widehat\nu_{u,v}(\omega)$ continuous. By (1),
$\nu_{0,1} = \mu_{\Sigma,\Sigma'}$ is infinitely divisible. Under (P) this holds for every pair, and
$\mu_\Sigma = \mu_{0,\Sigma}$ is the increment of the pair $0 \le_W \Sigma$.

\emph{(3)} is \ref{lem:path-condition-orders}, proved there.

\emph{(4).} For $s \le t$ put $\nu_{u,v} := \mu_{s + u(t-s),\,s + v(t-s)}$, $0 \le u \le v \le 1$.
(I1) and (I2) give the hypotheses of (1), and $\nu_{0,1} = \mu_{s,t}$.
\end{proof}

% CHANGED Q-0281: split from thm:cascade-infinitely-divisible (ADR-0025;
% records/formalization/PRICING.md): clause (3), the orders with the path condition and their chains,
% copied verbatim. It is [T] -- matrix analysis, no ledger entry -- where the parent is [A] on A6, and
% thm:cone-orders reads it. The parent's text and numbering are unchanged. Re-copied after the
% merge of main (Q-0287/Q-0289 split clause (c) into three sentences and moved the partial-order
% justification to the proof, without changing what it claims).
% CHANGED Q-0281: \uses re-pointed to the split child it reads (prop:stable-step-cone(2), the order, only).
% CHANGED Q-0281: blind review R1-2: the edge to the [A] parent is dropped (this [T] node read only
% the parent's notation, the W-chain of its clause (2), and the parent's proof text), and the W-chain
% is restated here instead of "use the notation of thm:cascade-infinitely-divisible(2)"; "satisfy the
% hypothesis of (2) for every pair" is written out as what it says (a partial order on the set named,
% every pair joined by a W-chain). The parent \uses this node. R1-6: (a)-(b)'s partial orders are
% lem:gw-scheme-claims(2), whose declarations the tag names. Nothing claimed changes.
\begin{lemma}[orders with the path condition]
\label{lem:path-condition-orders}
\uses{def:gw-cascade-family,lem:gw-scheme-claims,lem:stable-step-order}
\lean{AffineHemigroup.chain_loewner, AffineHemigroup.chain_ray_zero, AffineHemigroup.chain_ray,
AffineHemigroup.chain_stepSet, AffineHemigroup.stepSet_loewner, AffineHemigroup.stepSet_ray,
AffineHemigroup.stepSet_stable, AffineHemigroup.isWedgeDatum_loewner,
AffineHemigroup.isWedgeDatum_ray, AffineHemigroup.glSet_isClosedSubgroupGL}\leanok
Let $d \ge 1$. For a set $\mathcal P \subseteq \PSD(d)$ and a partial order $\le_W$ on $\mathcal P$,
call a continuous $\zeta : [0,1] \to \mathcal P$ with $\zeta(s) \le_W \zeta(t)$ for all
$s \le t$ a \emph{$W$-chain} from $\zeta(0)$ to $\zeta(1)$. Each of the following is a partial
order $\le_W$ on the set $\mathcal P$ named in which every pair $\Sigma \le_W \Sigma'$ is joined by a
$W$-chain, the chain being written out.
(a) The Loewner order on $\PSD(d)$: $\zeta(s) := \Sigma + s(\Sigma' - \Sigma)$.
(b) The ray order on $\PSD(d)$ (\ref{def:gw-cascade-family}): $\zeta(s) := s\,\Sigma'$ from
$\Sigma = 0$, and $\zeta(s) := (1 + s(c - 1))\,\Sigma$ from $\Sigma$ to $\Sigma' = c\,\Sigma$.
(c) Let $\mathcal K \subseteq \PD(d)$ satisfy $QMQ^{\tp} \in \mathcal K$ and $M^\tau \in \mathcal K$
whenever $M \in \mathcal K$, $Q \in \Orth(d)$ and $\tau \in [0,1]$. On $\mathcal P = \PD(d)$ write
$\Sigma \le_{\mathcal K} \Sigma'$ when $\Sigma^{-1/2}\Sigma'\Sigma^{-1/2} \in \mathcal K$, and suppose
that this relation is a partial order. Then every pair $\Sigma \le_{\mathcal K} \Sigma'$ is joined
by a $W$-chain for $\le_W := \le_{\mathcal K}$, namely
$\zeta(s) := \Sigma^{1/2}M^s\Sigma^{1/2}$ with $M := \Sigma^{-1/2}\Sigma'\Sigma^{-1/2}$.
This covers the Loewner order on $\PD(d)$ ($\mathcal K = \{M : M \ge I\}$), the ray order on $\PD(d)$
($\mathcal K = \{cI : c \ge 1\}$), and for $0 < \alpha < 2$ the order of the stable step cone,
$\mathcal K_\alpha := \{M \in \PD(d) : \lambda_{\min}(M)^{d+\alpha} \ge \det M\}$ with $\lambda_{\min}(M)$ the
least eigenvalue; $\mathcal K_\alpha$ is the set of $e^X$ with $X$ symmetric and
$(d+\alpha)\lambda_{\min}(X) \ge \tr X$. These three relations are partial orders.
\statusT\quad\emph{(Cites nothing beyond the spectral theorem. Clause (3) of
\ref{thm:cascade-infinitely-divisible}, split off on 2026-10-04 (Q-0281, ADR-0025): it cites no
source, where the parent's (1), (2) and (4) spend Khintchine's theorem; the numbering stays the
parent's.)}
\end{lemma}

% CHANGED Q-0281 wave 2: proved (AffineHemigroup.PathCondition; guarded, Lean core). The proof is
% rewritten to the Lean route: the relative step along the chain is an orthogonal conjugate of
% M^{t-s} by an explicit orthogonal matrix, so the conjugacy of isospectral symmetric matrices that
% the parent's proof cites is not used; the three step sets are read in the spectral form.
\begin{proof}
\leanok
Each $\zeta$ is continuous, with $\zeta(0) = \Sigma$ and $\zeta(1) = \Sigma'$, and takes values in
$\mathcal P$. That the Loewner and the ray order are partial orders on $\PSD(d)$ is
\ref{lem:gw-scheme-claims}(2).

(a) $\zeta(s) = (1-s)\Sigma + s\Sigma'$ is positive semidefinite for $s \in [0,1]$, and for
$s \le t$, $\zeta(t) - \zeta(s) = (t - s)(\Sigma' - \Sigma)$ is positive semidefinite.

(b) From $0$: $\zeta(0) = 0$ is below everything, and for $0 < s \le t$, $\zeta(t) = (t/s)\zeta(s)$.
From $\Sigma$ to $c\,\Sigma$ with $c \ge 1$: for $s \le t$, $\zeta(t) = \lambda\zeta(s)$ with
$\lambda = (1 + t(c-1))/(1 + s(c-1)) \ge 1$.

(c) Write $M = U\operatorname{diag}(m_1, \dots, m_d)U^{\tp}$ with $U$ orthogonal and $m_i > 0$
(the spectral theorem); then $M^s = U\operatorname{diag}(m_1^s, \dots, m_d^s)U^{\tp}$, which is
positive definite, continuous in $s$, symmetric, and satisfies $M^aM^b = M^{a+b}$, $M^0 = I$,
$M^1 = M$. So $\zeta$ is continuous with values in $\PD(d)$, $\zeta(0) = \Sigma$ and
$\zeta(1) = \Sigma^{1/2}M\Sigma^{1/2} = \Sigma'$. Let $0 \le s \le t \le 1$ and put
$B := \Sigma^{1/2}M^{s/2}$, so that $\zeta(s) = BB^{\tp}$ and
$\zeta(t) = \Sigma^{1/2}M^{s/2}M^{t-s}M^{s/2}\Sigma^{1/2} = BM^{t-s}B^{\tp}$. With
$Q := \zeta(s)^{-1/2}B$,
\[
  \zeta(s)^{-1/2}\zeta(t)\zeta(s)^{-1/2} = QM^{t-s}Q^{\tp},
  \qquad QQ^{\tp} = \zeta(s)^{-1/2}\zeta(s)\zeta(s)^{-1/2} = I,
\]
so $Q \in \Orth(d)$, and the relative step of $(\zeta(s), \zeta(t))$ lies in $\mathcal K$ by the
two closure properties with $\tau = t - s \in [0,1]$: $\zeta(s) \le_{\mathcal K} \zeta(t)$.

The three sets. For $V$ orthogonal and $w \in \RR^d$ write $D_V(w) := V\operatorname{diag}(w)V^{\tp}$;
then $\langle y, D_V(w)y\rangle = \sum_i w_i (V^{\tp}y)_i^2$ and $|V^{\tp}y| = |y|$, so
$\lambda_{\min}(D_V(w)) = \min_i w_i$, while $\det D_V(w) = \prod_i w_i$ and
$\tr D_V(w) = \sum_i w_i$. An orthogonal conjugate of $D_U(w)$ is $D_{QU}(w)$, with the same
least eigenvalue and determinant; this gives the closure of the three sets under orthogonal
conjugation ($QMQ^{\tp} - I = Q(M - I)Q^{\tp}$ for the first, $Q(cI)Q^{\tp} = cI$ for the second). For
$\{M : M \ge I\}$: $M - I = D_U(m - 1)$, so every $m_i \ge 1$, hence $m_i^\tau \ge 1$ and
$M^\tau \ge I$; and $\Sigma^{-1/2}\Sigma'\Sigma^{-1/2} - I = \Sigma^{-1/2}(\Sigma' - \Sigma)\Sigma^{-1/2}$, so
$\le_{\mathcal K}$ is the Loewner order on $\PD(d)$, a partial order. For $\{cI : c \ge 1\}$:
$D_U(m) = cI$ forces every $m_i = c$, so $(cI)^\tau = c^\tau I$ with $c^\tau \ge 1$; and
$\Sigma^{-1/2}\Sigma'\Sigma^{-1/2} = cI$ exactly when $\Sigma' = c\,\Sigma$, so $\le_{\mathcal K}$ is
the ray order on $\PD(d)$, a partial order. For $\mathcal K_\alpha$:
$\det M^\tau = (\det M)^\tau$ and $\lambda_{\min}(M^\tau) = \min_i m_i^\tau = \lambda_{\min}(M)^\tau$, so the inequality
defining $\mathcal K_\alpha$ is raised to the power $\tau \ge 0$; the relation is a partial order by
\ref{lem:stable-step-order}. Finally, $e^{D_U(x)} = D_U(e^{x_1}, \dots, e^{x_d})$, so a symmetric
$X = D_U(x)$ has $\lambda_{\min}(e^X) = e^{\lambda_{\min}(X)}$ and $\det e^X = e^{\tr X}$, and $e^X \in
\mathcal K_\alpha$ exactly when $(d+\alpha)\lambda_{\min}(X) \ge \tr X$; conversely $M = D_U(m) \in
\mathcal K_\alpha$ is $e^X$ with $X := D_U(\log m_1, \dots, \log m_d)$, for which
$\lambda_{\min}(X) = \log \lambda_{\min}(M)$ and $\tr X = \log\det M$.
\end{proof}

\section{The single step and rigidity}
\label{sec:affine-single-step}

The nodes are arranged so that the theorem comes last, as a corollary of a statement about a
\emph{single} step: an increment that widens the aperture in some directions and leaves the others
unchanged can be a probability measure only for the Gaussian. That statement is the draft's
Proposition~2.14 (\ref{thm:rank-deficient-step}), and its proof needs neither the isotropic
classification nor infinite divisibility; so neither is used.

% draft: section 2.2, the probability-measure form of Lemma 2.7, which is the form the draft's
% own proof of Proposition 2.14 uses. The draft's Lemma 2.7 proper (a negative definite psi
% vanishing on H) is not a node: nothing below needs it.
% CHANGED round-2 review (Gemini), item 15 (records/affine/reviews/VERIFY-round2-notation-paper.md):
% the frequency on R^d is omega, as the paper's rule asks from its section 3 on (was xi; and z, z_0
% for the zero in the proof of (1), now omega_1, omega_0). Notation only; no step changes.
\begin{lemma}[a characteristic function equal to $1$ on a subspace]
\label{lem:charfn-one-on-subspace}
\lean{AffineHemigroup.charFun_one_on_subspace,
AffineHemigroup.eq_dirac_of_charFun_one_near_zero}\leanok
Let $\rho$ be a probability measure on $\RR^d$ and $H \subseteq \RR^d$ a linear subspace with
$\widehat\rho(\eta) = 1$ for every $\eta$ in a neighbourhood of $0$ in $H$. Then $\rho$ is carried
by the orthogonal complement $H^\perp$, and $\widehat\rho(\omega + \eta) = \widehat\rho(\omega)$ for all
$\omega \in \RR^d$ and $\eta \in H$. In particular, if $\widehat\rho = 1$ near $0$ then
$\rho = \delta_0$.
\statusT\quad\emph{(Cites nothing. The hypothesis asks for $\widehat\rho = 1$ only near $0$ in
$H$, which \ref{prop:fixed-frame-families} needs, the transforms there possibly having zeros; every
other node reads it on all of $H$. The draft's Lemma~2.7 is the same statement for a
negative definite function $\psi$ and the \Levy{} measure of $e^{-\psi}$; it would need the
\Levy--Khintchine representation on $\RR^d$, which is \ledger{A3}'s territory. \textbf{Not
claimed}: that form. The probability-measure form is what the draft's proof of its
Proposition~2.14 already uses, and it is the only form any node here reads; so the draft's check
item~1(iii), the countable dense subset, is discharged here and needs no measure but $\rho$.)}
\end{lemma}

% CHANGED Q-0281 wave 1: proof rewritten to the Lean route (a basis of H and a countable family of
% its multiples in place of a countable dense subset of the ball and its connectedness).
\begin{proof}
\leanok
Let $r > 0$ be such that $\widehat\rho(\eta) = 1$ for $\eta \in H$ with $|\eta| < r$. For such an
$\eta$, $0 = \mathrm{Re}\bigl(1 - \widehat\rho(\eta)\bigr)
= \int\bigl(1 - \cos\langle\eta,x\rangle\bigr)\rho(dx)$ with a nonnegative integrand, so
$\cos\langle\eta,x\rangle = 1$ for every $x$ outside a $\rho$-null set $N_\eta$. Let
$v_1, \dots, v_k$ be a basis of $H$, put $c_j := r/(2(|v_j| + 1)) > 0$ and
$\eta_{j,m} := c_jv_j/(m+1)$ for $1 \le j \le k$ and integers $m \ge 0$; each $\eta_{j,m}$ lies
in $H$ and has $|\eta_{j,m}| < r$. The set $N := \bigcup_{j,m}N_{\eta_{j,m}}$ is a countable
union of null sets, hence $\rho$-null. Fix $x \notin N$ and $j$, and put
$\theta := c_j\langle v_j,x\rangle$; then $\cos(\theta/(m+1)) = 1$ for every $m \ge 0$. Choosing
$m$ with $|\theta|/(m+1) < 2\pi$, the only point of $(-2\pi, 2\pi)$ at which the cosine equals
$1$ is $0$, so $\theta = 0$, and $\langle v_j,x\rangle = 0$ as $c_j \ne 0$. As this holds for
every basis vector, $\langle\eta,x\rangle = 0$ for every $\eta \in H$, that is $x \in H^\perp$.
Thus $\rho(\RR^d \setminus H^\perp) \le \rho(N) = 0$, and for $\omega \in \RR^d$ and $\eta \in H$,
$\widehat\rho(\omega+\eta) = \int_{H^\perp}e^{i\langle\omega,x\rangle}e^{i\langle\eta,x\rangle}\rho(dx)
= \widehat\rho(\omega)$. For $H = \RR^d$ this gives $\widehat\rho(\eta) = \widehat\rho(0) = 1 =
\widehat{\delta_0}(\eta)$ for every $\eta$, so $\rho = \delta_0$ by the uniqueness of the
transform.
\end{proof}

% draft: section 2.2, Lemma 2.6, without its detour. The draft reaches (2.2) through Theorem 2.2 on
% the isotropic slice; covariance alone gives it, and the nonvanishing that makes the logarithm
% legitimate is the cascade plus continuity, as in [V, Lem. 4.2]. The draft's last sentence (the
% increments are infinitely divisible, so psi is negative definite) is not a clause: nothing needs
% it.
\begin{lemma}[reduction to one function]
\label{lem:affine-reduction}
\uses{def:affine-cascade-family,lem:charfn-one-on-subspace}
\lean{AffineHemigroup.affine_reduction}\leanok
Let $(\mu_{\Sigma,\Sigma'})$ be an affine cascade family in dimension $d$. There is a continuous
$g : [0,\infty) \to [0,\infty)$ with $g(0) = 0$ and $g \not\equiv 0$ such that
\begin{equation}\label{eq:affine-reduction}
  \widehat\mu_\Sigma(\omega) = e^{-g(\langle\omega,\Sigma\omega\rangle)}
  \qquad (\Sigma \in \PSD(d),\ \omega \in \RR^d),
\end{equation}
and for every pair $\Sigma \le \Sigma'$
\begin{equation}\label{eq:affine-increment}
  \widehat\mu_{\Sigma,\Sigma'}(\omega)
   = \exp\bigl[-g(\langle\omega,\Sigma'\omega\rangle) + g(\langle\omega,\Sigma\omega\rangle)\bigr] .
\end{equation}
\statusT\quad\emph{(Cites nothing. \textbf{Not claimed}: that $g(u) = F(\sqrt u)$ for an $F$ of
the isotropic classification's polar form (2.1), nor that the isotropic slice satisfies
\ref{def:isotropic-kernel-family} --- both go through \ref{thm:isotropic-classification} and its
reading \ref{prop:isotropic-polar-reading} in \S\ref{sec:isotropic-classification}, neither of
which this proof uses --- nor that the increments are infinitely divisible, the draft's last
sentence of Lemma~2.6, which is \ref{thm:cascade-infinitely-divisible}, by the null array in $\RR^d$
(ledger entry A6; the draft's check item~1(ii)). The proof below reads none of the three: the rigidity
below needs only that one increment is a probability measure.)}
\end{lemma}

% CHANGED Q-0281 wave 1: proof rewritten to the Lean route in its first paragraph (the reflection
% that carries one vector of a sphere to another names the orthogonal matrix; reality is the radial
% symmetry at -omega rather than a separate use of Q = -I). The rest is the Lean route as it stood.
\begin{proof}
\leanok
\emph{The kernel from scale zero at $I$ is radial and real.} For $Q \in \Orth(d)$, (K3) of
\ref{def:affine-cascade-family} with $\Sigma = 0$,
$\Sigma' = I$ gives $Q_*\mu_I = \mu_{QQ^{\tp}} = \mu_I$, so
$\widehat\mu_I(Q^{\tp}\omega) = \widehat\mu_I(\omega)$. Given $\omega, \omega'$ with
$|\omega| = |\omega'|$, the reflection $R$ in the hyperplane $(\omega - \omega')^\perp$ is
orthogonal and symmetric and maps $\omega$ to $\omega'$; with $Q = R$ this gives
$\widehat\mu_I(\omega') = \widehat\mu_I(\omega)$, so $\widehat\mu_I$ is constant on spheres.
Put $\vartheta(u) := \widehat\mu_I(\sqrt u\,e_1)$ for $u \ge 0$; then
$\widehat\mu_I(\omega) = \vartheta(|\omega|^2)$, $\vartheta$ is continuous because a characteristic
function is, and $\vartheta(0) = 1$. As $|-\omega| = |\omega|$,
$\widehat\mu_I(\omega) = \widehat\mu_I(-\omega) = \overline{\widehat\mu_I(\omega)}$, so $\vartheta$ is
real, and $|\vartheta| \le 1$.

\emph{Every kernel from scale zero.} For $\Sigma \in \PD(d)$, (K3) with $A = \Sigma^{1/2}$ gives
$\mu_\Sigma = A_*\mu_I$, so $\widehat\mu_\Sigma(\omega) = \widehat\mu_I(\Sigma^{1/2}\omega)
= \vartheta(\langle\omega,\Sigma\omega\rangle)$. For $\Sigma \in \PSD(d)$ singular,
$\Sigma_n := \Sigma + n^{-1}I \in \PD(d)$ tends to $\Sigma$, and (K2) along the pairs
$(0,\Sigma_n) \to (0,\Sigma)$ with the continuity of $\vartheta$ gives
$\widehat\mu_\Sigma(\omega) = \lim_n\vartheta(\langle\omega,\Sigma_n\omega\rangle)
= \vartheta(\langle\omega,\Sigma\omega\rangle)$.

\emph{No zeros.} Fix $\omega$ and put $f(s,t) := \widehat\mu_{sI,tI}(\omega)$ for $0 \le s \le t$.
By (K2) $f$ is continuous; by (K1), and because the transform of a convolution is the product of
the transforms, $f(t,t) = 1$ and $f(r,t) = f(s,t)\,f(r,s)$ for $r \le s \le t$. Suppose
$f(0,t) = 0$ for some $t$, and let $t^*$ be the infimum of such $t$. As $f(0,\cdot)$ is continuous
with $f(0,0) = 1$, $t^* > 0$ and $f(0,t^*) = 0$. By continuity at $(t^*,t^*)$, where $f = 1$, there
is $s \in (0,t^*)$ with $f(s,t^*) \ne 0$, and $f(0,s) \ne 0$ because $s < t^*$; then
$f(0,t^*) = f(s,t^*)\,f(0,s) \ne 0$, a contradiction. Hence
$\vartheta(t|\omega|^2) = \widehat\mu_{tI}(\omega) \ne 0$ for all $t \ge 0$ and $\omega$, so $\vartheta$
has no zero on $[0,\infty)$. Being real and continuous with $\vartheta(0) = 1$, it is positive, and
$g := -\log\vartheta$ is continuous, nonnegative because $\vartheta \le 1$, and $g(0) = 0$. This is
\ref{eq:affine-reduction}.

\emph{Increments.} For $\Sigma \le \Sigma'$, (K1) gives $\mu_{\Sigma'} = \mu_{\Sigma,\Sigma'} *
\mu_\Sigma$, so $\widehat\mu_{\Sigma'} = \widehat\mu_{\Sigma,\Sigma'}\,\widehat\mu_\Sigma$, and
$\widehat\mu_\Sigma > 0$ gives \ref{eq:affine-increment}.

\emph{Nondegeneracy.} If $g \equiv 0$ then $\widehat\mu_I \equiv 1$ and
\ref{lem:charfn-one-on-subspace} with $H = \RR^d$ gives $\mu_I = \delta_0$, against (K4).
\end{proof}

% draft: section 2.2, the functional-equation step of Theorem 2.8's proof and of Proposition
% 2.14's, made one lemma. The single-ratio form is Proposition 2.14's (2.5); the whole-order form
% (2.4) is the special case in which every a > 1 is available.
\begin{lemma}[one ratio forces linearity]
\label{lem:single-ratio-cauchy}
\lean{AffineHemigroup.single_ratio_cauchy}\leanok
Let $g : [0,\infty) \to [0,\infty)$ be continuous at $0$ with $g(0) = 0$, and let $a > 1$ be such
that
\begin{equation}\label{eq:single-ratio}
  g(au + w) - g(u + w) = g(au) - g(u) \qquad \text{for all } u > 0,\ w \ge 0 .
\end{equation}
Then $g(u) = b\,u$ for every $u \ge 0$, with $b = g(1) \ge 0$.
\statusT\quad\emph{(Cites nothing. The hypotheses are the draft's: continuity at $0$ only, and
nonnegativity, which replaces every other regularity in Cauchy's equation. The draft's check
items~1(iv) and~11 ask whether continuity at $0$ suffices to identify the constant with
$g(\delta)$ and whether the telescoping from a single ratio is sound; the proof answers both, the
first by halving rather than by a limit at $\delta$.)}
\end{lemma}

% CHANGED Q-0281 wave 1: proof rewritten to the Lean route in its last two steps (the halving through
% h = g - K, and the squeeze through floor(nx) in place of rationals and a limit).
\begin{proof}
\leanok
\emph{A constant difference on a half-line.} Fix $\delta > 0$ and put $u_\delta := \delta/(a-1)$,
so that $a u_\delta = u_\delta + \delta$. For $x \ge u_\delta$, \ref{eq:single-ratio} at
$u = u_\delta$ and $w = x - u_\delta \ge 0$ reads
$g(x+\delta) - g(x) = g(u_\delta + \delta) - g(u_\delta) =: \kappa(\delta)$, independent of $x$ on
$[u_\delta,\infty)$.

\emph{A constant difference everywhere.} Fix $\delta > 0$ and $y > 0$. For every integer
$n \ge \delta/((a-1)y)$ the points $y + j\delta/n$, $0 \le j < n$, are at least
$y \ge u_{\delta/n}$, so the previous step with $\delta/n$ in place of $\delta$ gives
\[
  g(y+\delta) - g(y) = \sum_{j=0}^{n-1}\Bigl[g\bigl(y + \tfrac{(j+1)\delta}{n}\bigr)
     - g\bigl(y + \tfrac{j\delta}{n}\bigr)\Bigr] = n\,\kappa(\delta/n) .
\]
Given $y_1, y_2 > 0$, an $n$ admissible for both shows
$g(y_1+\delta) - g(y_1) = g(y_2+\delta) - g(y_2)$. So $K(\delta) := g(y+\delta) - g(y)$ does not
depend on $y > 0$, and $K(2\delta) = [g(y+2\delta) - g(y+\delta)] + [g(y+\delta) - g(y)]
= 2K(\delta)$.

\emph{Additivity, from continuity at $0$ alone.} Put $h(\delta) := g(\delta) - K(\delta)$ for
$\delta > 0$. With $y = \delta$ the previous step gives $g(2\delta) = g(\delta) + K(\delta)$, so
$h(2\delta) = g(\delta) + K(\delta) - 2K(\delta) = h(\delta)$. By induction, for every $n \ge 0$,
\[
  h(\delta) = h(\delta 2^{-n}) , \qquad K(\delta) = 2^nK(\delta 2^{-n}) , \qquad\text{hence}\qquad
  h(\delta) = g(\delta 2^{-n}) - 2^{-n}K(\delta) .
\] Letting $n \to \infty$, continuity at $0$ and
$g(0) = 0$ give $h(\delta) = 0$, that is $g(\delta) = K(\delta)$. Hence
$g(y + \delta) = g(y) + g(\delta)$ for all $y, \delta > 0$.

\emph{Cauchy's equation.} As $g \ge 0$, $g(y+\delta) \ge g(y)$: $g$ is nondecreasing on
$[0,\infty)$ (at $0$ because $g(0) = 0 \le g$). By induction $g(m) = m\,g(1)$ for integers
$m \ge 0$ and $g(nx) = n\,g(x)$ for $x > 0$ and integers $n \ge 1$. Fix $x > 0$ and $n \ge 1$, and
let $m$ be the largest integer with $m \le nx$, so that $m \le nx < m + 1$. Monotonicity gives
$m\,g(1) \le n\,g(x) \le (m+1)\,g(1)$, and $g(1) \ge 0$ gives
$m\,g(1) \le nx\,g(1) \le (m+1)\,g(1)$; so $n\,|g(x) - x\,g(1)| \le g(1)$ for every $n \ge 1$, which
forces $g(x) = x\,g(1)$. With $g(0) = 0$ that is the claim.
\end{proof}

% draft: section 2.4, Proposition 2.14, moved forward: it is the engine of Theorem 2.8. The
% hypotheses are weakened: the draft asks that g come from an admissible F, and the proof reads
% only g >= 0, continuity at 0 and one factorization; exp(-g(<w,Sigma w>)) need not even be a
% characteristic function.
\begin{theorem}[one rank-deficient step forces the Gaussian]
\label{thm:rank-deficient-step}
\uses{lem:charfn-one-on-subspace,lem:single-ratio-cauchy}
\lean{AffineHemigroup.rank_deficient_step,
AffineHemigroup.rank_deficient_step_converse}\leanok
Let $d \ge 2$ and let $g : [0,\infty) \to [0,\infty)$ be continuous at $0$ with $g(0) = 0$.
Suppose there are $\Sigma \in \PD(d)$, $P \in \PSD(d)$ with $0 < \rank P < d$, and a probability
measure $\rho$ on $\RR^d$ with
\begin{equation}\label{eq:one-step}
  e^{-g(\langle\omega,(\Sigma+P)\omega\rangle)}
   = e^{-g(\langle\omega,\Sigma\omega\rangle)}\,\widehat\rho(\omega)
   \qquad (\omega \in \RR^d).
\end{equation}
Then $g(u) = b\,u$ with $b = g(1) \ge 0$, and $\rho$ is the centred Gaussian law with covariance
$2bP$. Conversely, if $g(u) = bu$ with $b \ge 0$, then \ref{eq:one-step} holds with that Gaussian
$\rho$ for every $\Sigma, P \in \PSD(d)$.
\statusT\quad\emph{(Cites nothing beyond the uniqueness of the Fourier transform of a
probability measure on $\RR^d$, used to name $\rho$. No hemigroup, no continuity in $\Sigma$ and
no infinite divisibility of $\rho$ is assumed, as the draft says; nor, here, that $g$ comes from
an admissible exponent. \textbf{Not claimed}: anything about steps of full rank, where every
directional variance grows and $\ker P = \{0\}$ leaves the subspace argument nothing to act on.
For the \Matern{} orbit of index $1$ those fail too, by a different argument,
\ref{prop:matern-cross-shape}; for the isotropic $\alpha$-stable orbit a cone of them succeeds,
\ref{prop:stable-step-cone}, which refutes the draft's conjecture at the end of its \S 2.4 that
they fail for every non-Gaussian $g$. So the full-rank case depends on $g$.)}
\end{theorem}

% CHANGED Q-0281 wave 1: proof rewritten to the Lean route (two vectors chosen directly in place of
% the transport to the identity and the diagonalization of Sigma^{-1/2} P Sigma^{-1/2}).
\begin{proof}
\leanok
Write $q_\Sigma(\omega) := \langle\omega,\Sigma\omega\rangle$ and
$q_P(\omega) := \langle\omega,P\omega\rangle$, so that $q_{\Sigma+P} = q_\Sigma + q_P$.

\emph{The transform of $\rho$.} Dividing \ref{eq:one-step} by
$e^{-g(q_\Sigma(\omega))}$, which has no zero,
\[
  \widehat\rho(\omega) = \exp\bigl[-g\bigl(q_\Sigma(\omega) + q_P(\omega)\bigr)
     + g\bigl(q_\Sigma(\omega)\bigr)\bigr] \qquad (\omega \in \RR^d).
\]

\emph{Two vectors.} As $\rank P < d$, there is $\eta \ne 0$ with $P\eta = 0$. There is then a
$\xi$ with $\langle\xi,\Sigma\eta\rangle = 0$ and $P\xi \ne 0$: otherwise $P$ vanishes on the
hyperplane $(\Sigma\eta)^\perp$, and every $\omega$ is
$\omega = (\omega - c\eta) + c\eta$ with $c := \langle\omega,\Sigma\eta\rangle/q_\Sigma(\eta)$
(note $q_\Sigma(\eta) > 0$, $\Sigma$ being positive definite) and $\omega - c\eta$ in that
hyperplane, so $P\omega = 0$ for every $\omega$ and $P = 0$, against $\rank P > 0$. Then
$\xi \ne 0$ (as $P\xi \ne 0$), so $q_\Sigma(\xi) > 0$, and $q_P(\xi) = |P^{1/2}\xi|^2 > 0$, since
$P^{1/2}\xi = 0$ would give $P\xi = 0$; put $p := q_P(\xi)/q_\Sigma(\xi) > 0$ and $a := 1 + p > 1$.

\emph{The subspace.} On $H := \ker P$ one has $q_P = 0$, so the exponent vanishes and
$\widehat\rho = 1$ on $H$. By \ref{lem:charfn-one-on-subspace},
$\widehat\rho(\zeta + \eta') = \widehat\rho(\zeta)$ for $\zeta \in \RR^d$ and $\eta' \in H$. Given
$u > 0$ and $w \ge 0$, take $\zeta := t\xi$ and $\eta' := s\eta$ with
$t := (u/q_\Sigma(\xi))^{1/2}$ and $s := (w/q_\Sigma(\eta))^{1/2}$. Because
$\langle\xi,\Sigma\eta\rangle = 0$ and $P\eta = 0$, the quadratic forms split:
$q_\Sigma(t\xi + s\eta) = t^2q_\Sigma(\xi) + s^2q_\Sigma(\eta) = u + w$,
$q_P(t\xi + s\eta) = t^2q_P(\xi) = pu$, $q_\Sigma(t\xi) = u$ and $q_P(t\xi) = pu$. So the
invariance $\widehat\rho(t\xi + s\eta) = \widehat\rho(t\xi)$ reads
$e^{-g(u + w + pu) + g(u + w)} = e^{-g(u + pu) + g(u)}$, that is
\[
  g(au + w) - g(u + w) = g(au) - g(u) \qquad (u > 0,\ w \ge 0).
\]
That is \ref{eq:single-ratio}, and \ref{lem:single-ratio-cauchy} gives $g(u) = bu$ with
$b = g(1) \ge 0$.

\emph{The increment, and the converse.} Now \ref{eq:one-step} reads
$\widehat\rho(\omega) = e^{-b\langle\omega,P\omega\rangle}$, the characteristic function of the
centred Gaussian law with covariance $2bP$ (the law of $(2bP)^{1/2}Z$, $Z$ standard normal on
$\RR^d$); by uniqueness of the transform $\rho$ is that law. Conversely, for $g(u) = bu$ the
exponents in \ref{eq:one-step} are linear in the matrix, and the same Gaussian satisfies it for any
$\Sigma, P \in \PSD(d)$.
\end{proof}

% draft: section 2.2, Theorem 2.8, as a corollary of the single-step theorem. The draft's
% sentence "no isotropic family with a catalogue ... extends to an affine cascade family" is the
% last clause, stated on the kernel from scale zero so that it does not read the draft's
% Theorem 2.2.
% CHANGED Q-0281: \uses re-pointed to the split child it reads (the transform eq:affine-gaussian).
\begin{theorem}[affine rigidity: the affine Gaussian scale space]
\label{thm:affine-rigidity}
\uses{def:affine-cascade-family,lem:affine-reduction,thm:rank-deficient-step,lem:affine-gaussian-transform}
\lean{AffineHemigroup.affine_rigidity, AffineHemigroup.affine_rigidity_kernel}\leanok
Let $d \ge 2$. A family $(\mu_{\Sigma,\Sigma'})$ is an affine cascade family in dimension $d$ if
and only if it is the affine Gaussian scale space with some constant $b > 0$
(\ref{eq:affine-gaussian}). In particular the kernel from scale zero at $I$ of an affine cascade
family in dimension $d \ge 2$ is a centred Gaussian law with covariance a positive multiple of $I$,
and no other probability law on $\RR^d$ is the kernel $\mu_I$ of one.
\statusT\quad\emph{(Cites nothing beyond the uniqueness of the Fourier transform. The proof reads
(K1)--(K4) through \ref{lem:affine-reduction} and then \textbf{one} step,
$I \le I + e_1e_1^{\tp}$, of the cascade: as the draft's first comment says, no regularity of $g$
beyond continuity at $0$, no locality, no semigroup and no extremum principle is used, and here no
infinite divisibility either. \textbf{Not claimed here}: the theorem on operators, both halves,
which is \ref{thm:affine-rigidity-operators} (this theorem with \ref{lem:operators-to-kernels}(4)
and \ref{lem:kernels-to-operators}(2)); the
comparison with Lindeberg's derivation of the affine Gaussian scale space in the draft's
Corollary~2.9 --- Lindeberg's route assumes a semigroup and non-enhancement of local extrema,
neither used here, and read at the page (records/affine/READING-lindeberg2011generalized.md) the
comparison is accurate, but it stays prose, since no node's proof depends on it; the case $d = 1$, where the
hyperplane of the proof is $\{0\}$ and the draft expects the whole self-decomposable class of
[V, Thm.~7.3] to survive; and the draft's remark that a rank-one base with a step in a second
direction already suffices, which is not needed for the statement. \textbf{What it is worth.} The
proof is on kernels and reads (K1)--(K4) and one cascade step; it does not use
\ref{thm:isotropic-classification}, the node the isotropic classification in dimension $d$ now
has. Read against that node, this theorem says that no isotropic family with nonzero profile ---
no kernel from scale zero of an isotropic cascade family with $F \not\equiv 0$ in
\ref{thm:isotropic-classification}'s sense, hence none of the corners and none of the extreme rays
of \ref{prop:isotropic-bridge} --- extends to an affine cascade family, since its kernel $\mu_I$ is
not Gaussian. \textbf{Prior art, and a covariance-free form} (a remark, by the author's decision of
2026-10-03: no node and no ledger entry). The core is Kelker's characterization of the normal law.
From the base $I - e_1e_1^{\tp}$, $\mu_I = \mu_{I - e_1e_1^{\tp}} * \mu_{I - e_1e_1^{\tp},\,I}$ with
the first factor carried by $e_1^\perp$ and the second by the $e_1$-axis
(\ref{lem:charfn-one-on-subspace}), so the first coordinate of $\mu_I$ is independent of the
others; by rotation invariance so is every coordinate, and a spherical law with independent
coordinates is normal (Kelker 1970, Lemma~5, p.~423, read at the page; restated as Fang--Kotz--Ng
1990, Thm.~4.11, p.~106). Without any covariance group a weaker rigidity holds: if a kernel is
reached from scale zero along two decompositions into rank-one steps whose increments are
one-directional blurs, it is a sum of independent one-dimensional components in two frames, and
the Darmois--Skitovich theorem (Skitovich 1954, Thm.~2, p.~190, with Thm.~1, p.~185, for the case
in which every component is mixed; Comon 1994, Thm.~10, p.~294) makes Gaussian every non-degenerate
component that the two frames mix; the exceptional pairs of step directions are
those conjugate with respect to $\Sigma^{-1}$, which is this module's derivation and is stated by
neither source (records/affine/READING-skitovich1954linear.md, READING-comon1994independent.md).
Neither statement is used by a proof here.)}
\end{theorem}

\begin{proof}
\leanok
\emph{Necessity.} By \ref{lem:affine-reduction} there is a continuous $g \ge 0$ with $g(0) = 0$
and $g \not\equiv 0$ satisfying \ref{eq:affine-reduction} and \ref{eq:affine-increment}. Take
$\Sigma = I$ and $P = e_1e_1^{\tp}$, of rank $1 < d$, and $\rho := \mu_{I,I+P}$, a probability
measure. Then \ref{eq:affine-increment} at the pair $I \le I + P$, multiplied by
$e^{-g(|\omega|^2)}$, is \ref{eq:one-step}. So \ref{thm:rank-deficient-step} gives $g(u) = bu$,
and $b > 0$ because $g \not\equiv 0$. Now \ref{eq:affine-increment} is \ref{eq:affine-gaussian} for
every pair, and by uniqueness of the transform each $\mu_{\Sigma,\Sigma'}$ is the stated Gaussian.

\emph{Sufficiency.} Let $b > 0$ and let $\mu_{\Sigma,\Sigma'}$ be the centred Gaussian with
covariance $2b(\Sigma'-\Sigma)$, which exists for every positive semidefinite covariance as the law
of $(2b(\Sigma'-\Sigma))^{1/2}Z$. Its transform is \ref{eq:affine-gaussian}. The axioms of
\ref{def:affine-cascade-family} hold. (K1): the transform
at $\Sigma = \Sigma'$ is $1$, and the exponents of $\mu_{\Sigma',\Sigma''}$ and
$\mu_{\Sigma,\Sigma'}$ add to that of $\mu_{\Sigma,\Sigma''}$; by uniqueness of the transform,
$\mu_{\Sigma,\Sigma} = \delta_0$ and the convolution identity holds. (K2): the transform is
continuous in $(\Sigma,\Sigma')$. (K3): $\widehat{A_*\mu_{\Sigma,\Sigma'}}(\omega)
= e^{-b\langle A^{\tp}\omega,(\Sigma'-\Sigma)A^{\tp}\omega\rangle}
= e^{-b\langle\omega,(A\Sigma'A^{\tp} - A\Sigma A^{\tp})\omega\rangle}$. (K4):
$\widehat\mu_{0,I}(e_1) = e^{-b} \ne 1$.

\emph{The last clause.} By necessity, $\widehat\mu_I(\omega) = e^{-b|\omega|^2}$ with $b > 0$.
\end{proof}

% not in the draft: round 3 of the external review, report B, part 3 (ii) and its required change 2
% (records/affine/reviews/review-content-review-3-B-section-3-appendix-A-claude-fable-5.1.md;
% VERIFY-round3-front-classical-paper.md, B.S2), the author's decision D12 (PLAN-review-response-paper.md
% R3.0, 2026-11-03): where continuity (K2) enters thm:affine-rigidity. Asked as a remark; it is a node,
% with a proof, because its paragraphs make two claims and prove them, and the blueprint has no remark
% environment. What writing it down found: (K2) is read in the proof of thm:affine-rigidity twice, in
% lem:affine-reduction, to extend the reduction to the singular apertures and to exclude a first zero
% along the ray of I; the second use is replaced by the self-decomposability of mu_I, which (K1) and (K3)
% give at the dilations (the argument of lem:gw-kernels-self-decomposable), so only the first remains.
% The family of (2) is the referee's, checked here by hand and numerically
% (records/affine/reviews/verify-round3-paper-scripts/k2_boundary_family.py: d = 2, 3, 12 000 Loewner
% pairs each). Priced at writing: clause (1) is an assembly of the Lean of lem:affine-reduction (its
% first two steps on PD(d)), gw_kernels_selfDecomposable, selfDecomposable_charFun_ne_zero and
% rank_deficient_step, about 80-120 lines, Lean core; clause (2) is Gaussian arithmetic, about 60-100
% lines, Lean core. No declaration is tagged until one is written (\notready, proved in print).
% CHANGED round 3, the Lean: proved, Lean core, on the proof of record's route, the "No declaration is
% tagged" above superseded. AffineHemigroup.RigidityBoundary, 505 lines against the 140-220 priced: 123 for
% clause (1) with its two lemmas, 281 for clause (2), whose (K1), (K3), (K4), the failure of the scale
% space and the failure of (K2) are a lemma each, and 101 for the header, the hypothesis structure and the
% helpers. The hypothesis is IsAffineCascadeFamilyNoK2, IsGWCascadeFamily without continuous at glSet d
% and loewnerLE. The text is unchanged.
\begin{proposition}[continuity is needed only at the boundary of $\PSD(d)$]
\label{prop:affine-rigidity-boundary}
\uses{def:affine-cascade-family,def:isotropic-kernel-family,lem:affine-reduction,prop:gw-ray-families,lem:charfn-one-on-subspace,thm:rank-deficient-step}
\lean{AffineHemigroup.affine_rigidity_interior, AffineHemigroup.affine_rigidity_boundary_family,
  AffineHemigroup.IsAffineCascadeFamilyNoK2, AffineHemigroup.boundaryFamily}\leanok
Let $d \ge 2$, and let $(\mu_{\Sigma,\Sigma'})$ be probability measures on $\RR^d$, indexed by the pairs
$\Sigma \le \Sigma'$ in $\PSD(d)$, that satisfy (K1), (K3) and (K4) of \ref{def:affine-cascade-family}, but
not necessarily (K2).
\begin{enumerate}
\item \emph{(Between positive definite apertures.)} There is $b > 0$ such that for every $\Sigma \in \PD(d)$
the kernel from scale zero $\mu_\Sigma$ is the centred Gaussian law with covariance $2b\Sigma$, and for every
pair $\Sigma \le \Sigma'$ in $\PD(d)$ the increment $\mu_{\Sigma,\Sigma'}$ is the centred Gaussian law with
covariance $2b(\Sigma' - \Sigma)$. In particular $\mu_I$ is the centred Gaussian law with covariance $2bI$.
\item \emph{(At the boundary.)} Let $b > 0$ and $0 \le \theta_0 < 1$, put $\theta(\Sigma) := 1$ for
$\Sigma \in \PD(d)$ and $\theta(\Sigma) := \theta_0$ for singular $\Sigma \in \PSD(d)$, and for
$\Sigma \le \Sigma'$ let $\mu_{\Sigma,\Sigma'}$ be the centred Gaussian law with covariance
$2b(\theta(\Sigma')\Sigma' - \theta(\Sigma)\Sigma)$. This covariance is positive semidefinite for every pair,
and the family satisfies (K1), (K3) and (K4) and agrees with the affine Gaussian scale space with constant
$b$ on the pairs in $\PD(d)$. Its kernel from scale zero at a nonzero singular $P$ has covariance
$2b\theta_0P$, so the family is not an affine Gaussian scale space, and it does not satisfy (K2).
\end{enumerate}
\statusT\quad\emph{(Cites nothing beyond the uniqueness of the Fourier transform and the transform of a
centred Gaussian law. \textbf{What it is worth.} \ref{thm:affine-rigidity} assumes continuity in the pair,
and Iijima's axioms, as Weickert, Ishikawa and Imiya report them (\sscite{weickert1999linear}, p.~241), assume
none. Clause (1) says that on
the interior $\PD(d)$ the scheme needs none either: the cascade, covariance and nondegeneracy force the
Gaussian between definite apertures, the self-decomposability of $\mu_I$ doing what continuity does in
\ref{lem:affine-reduction}. Clause (2) says that (K2) cannot be dropped at the boundary, whose kernels the
other axioms leave free, as the factor $\theta_0$ shows (with $\theta_0 = 0$ they are $\delta_0$). So (K2)
enters \ref{thm:affine-rigidity} only to pin the kernels and increments at singular apertures.
\textbf{Not claimed}: anything for a smaller group or a smaller order, where clause (1)'s single
rank-deficient step $I \le I + e_1e_1^{\tp}$ may be absent; and a description of all families satisfying
(K1), (K3) and (K4), which at the boundary may depend on more than one constant.)}
\end{proposition}

\begin{proof}
\leanok
\emph{(1).} Let $\vartheta(u) := \widehat\mu_I(\sqrt u\,e_1)$, $u \ge 0$. As in the first step of the proof of
\ref{lem:affine-reduction}, which reads only (K3), $\mu_I$ is rotation invariant,
$\widehat\mu_I(\omega) = \vartheta(|\omega|^2)$, and $\vartheta$ is continuous and real with $\vartheta(0) = 1$
and $|\vartheta| \le 1$. For $\Sigma \in \PD(d)$, (K3) with $A = \Sigma^{1/2}$ gives
$\widehat\mu_\Sigma(\omega) = \vartheta(\langle\omega,\Sigma\omega\rangle)$. For $0 < c < 1$, (K1) at
$0 \le c^2I \le I$ gives $\mu_I = \mu_{c^2I} * \mu_{c^2I,I}$, and (K3) with $A = cI$ at the pair $(0, I)$
gives $\mu_{c^2I} = (cI)_*\mu_I$. So $\mu_I$ is self-decomposable \ref{eq:self-decomposable}, and its
transform has no zero (\ref{prop:gw-ray-families}(1)). Being real and continuous with $\vartheta(0) = 1$,
$\vartheta$ is then positive, and $g := -\log\vartheta$ is continuous, nonnegative and $g(0) = 0$. If
$g \equiv 0$, then $\widehat\mu_I \equiv 1$ and $\mu_I = \delta_0$ (\ref{lem:charfn-one-on-subspace}), against
(K4). Put $P := e_1e_1^{\tp}$, of rank $1 < d$. (K1) at $0 \le I \le I + P$ gives
$\mu_{I+P} = \mu_{I,I+P} * \mu_I$, which is \ref{eq:one-step} with $\Sigma = I$ and $\rho = \mu_{I,I+P}$. So
\ref{thm:rank-deficient-step} gives $g(u) = bu$ with $b = g(1)$, and $b > 0$ as $g \not\equiv 0$. Hence
$\widehat\mu_\Sigma(\omega) = e^{-b\langle\omega,\Sigma\omega\rangle}$ for $\Sigma \in \PD(d)$, the transform of
the centred Gaussian law with covariance $2b\Sigma$. For $\Sigma \le \Sigma'$ in $\PD(d)$, (K1) at
$0 \le \Sigma \le \Sigma'$ gives $\widehat\mu_{\Sigma'} = \widehat\mu_{\Sigma,\Sigma'}\,\widehat\mu_\Sigma$, and as
$\widehat\mu_\Sigma$ has no zero, $\widehat\mu_{\Sigma,\Sigma'}(\omega) = e^{-b\langle\omega,(\Sigma'-\Sigma)\omega\rangle}$.
The uniqueness of the transform names each law.

\emph{(2), the covariances.} Let $\Sigma \le \Sigma'$ and $D := \theta(\Sigma')\Sigma' - \theta(\Sigma)\Sigma$. If
$\Sigma$ is definite, so is $\Sigma' \ge \Sigma$, and $D = \Sigma' - \Sigma \ge 0$. If both are singular,
$D = \theta_0(\Sigma' - \Sigma) \ge 0$. If $\Sigma$ is singular and $\Sigma'$ definite,
$D = \Sigma' - \theta_0\Sigma = (\Sigma' - \Sigma) + (1 - \theta_0)\Sigma \ge 0$. So each $\mu_{\Sigma,\Sigma'}$ exists,
with transform $e^{-b\langle\omega,D\omega\rangle}$.

\emph{(2), the axioms.} (K1): at $\Sigma = \Sigma'$, $D = 0$. For $\Sigma \le \Sigma' \le \Sigma''$ the
matrices $D$ of the pairs $(\Sigma, \Sigma')$ and $(\Sigma', \Sigma'')$ add to that of $(\Sigma, \Sigma'')$, so the
transforms multiply and the uniqueness of the transform gives the cascade identity. (K3): an $A \in \GL(d)$
preserves rank, so $\theta(A\Sigma A^{\tp}) = \theta(\Sigma)$, and $A_*\mu_{\Sigma,\Sigma'}$, the centred Gaussian
law with covariance $2bADA^{\tp}$, is $\mu_{A\Sigma A^{\tp},\,A\Sigma'A^{\tp}}$. (K4): $\theta(0)\,0 = 0$, so
$\mu_{0,I}$ has covariance $2bI$ and is not $\delta_0$. On a pair in $\PD(d)$, $D = \Sigma' - \Sigma$, which is
\ref{eq:affine-gaussian}. For a nonzero singular $P$, $\mu_{0,P}$ has covariance $2b\theta_0P$. An affine Gaussian scale space
with this $\mu_{0,I}$ has the constant $b$, by \ref{eq:affine-gaussian} at the pair $(0, I)$, and its kernel at $P$
has covariance $2bP \ne 2b\theta_0P$, as $\theta_0 < 1$ and $P \ne 0$. Finally, $P + n^{-1}I$ is
definite and tends to $P$, while for $\omega$ with $P\omega \ne 0$,
$\widehat\mu_{0,P+n^{-1}I}(\omega) = e^{-b\langle\omega,(P+n^{-1}I)\omega\rangle} \to e^{-b\langle\omega,P\omega\rangle}
\ne e^{-b\theta_0\langle\omega,P\omega\rangle} = \widehat\mu_{0,P}(\omega)$. So (K2) fails at the pair $(0, P)$.
\end{proof}

% not in the draft: queue item Q-0314 (hub), brief item 2: the operator form of the rigidity, both
% halves, which thm:affine-rigidity's annotation leaves unclaimed ("the sufficiency half on
% operators ... goes through the passage def:affine-cascade-family's annotation declines to claim").
% That passage is now lem:kernels-to-operators(2); the necessity half was lem:operators-to-kernels(4)
% followed by thm:affine-rigidity, said in the article's prose and not a node, and is now this node's
% first half. Added as a new node so that thm:affine-rigidity is unchanged while the Lean skeleton is
% typed (Q-0281). Cost, priced when the statement was written: assembly only, once
% lem:kernels-to-operators and lem:operators-to-kernels-families are declarations, about 40-80 Lean
% lines. No \lean tag.
% CHANGED Q-0281 wave M, phase 1: typed (the "No \lean tag" above is superseded). One declaration,
% AffineHemigroup.affine_rigidity_operators in AffineSkeleton/Ch04Classical.lean, covers the node:
% the equivalence, with "for every pair" read on the Loewner pairs P_W (the blind referee's note 1),
% and the uniqueness of b as a second conjunct, for any family represented by two affine Gaussian
% scale spaces, which is what the proof's appeal to Phi_{0,I} gives. Priced at typing: about
% 40-100 lines of assembly. The text is unchanged.
% CHANGED Q-0281 wave M, phase 2: proved, Lean core, on the proof of record's route
% (AffineHemigroup.RigidityOperators). The text is unchanged.
\begin{theorem}[affine rigidity on operators]
\label{thm:affine-rigidity-operators}
\uses{thm:affine-rigidity,def:gw-cascade-family,def:affine-cascade-family,lem:operators-to-kernels,lem:operators-to-kernels-families,lem:kernels-to-operators,lem:gw-scheme-claims,lem:affine-gaussian-transform}
\lean{AffineHemigroup.affine_rigidity_operators}\leanok
Let $d \ge 2$. A family $(\Phi_{\Sigma,\Sigma'})$ of linear operators on $L^1(\RR^d)$, indexed by the
pairs $\Sigma \le \Sigma'$ in $\PSD(d)$, is a $(\GL(d), \text{Loewner})$-cascade family of operators
(\ref{def:gw-cascade-family}(1)) if and only if there is $b > 0$ such that
$\Phi_{\Sigma,\Sigma'}f = \mu_{\Sigma,\Sigma'} * f$ for every pair and every $f \in L^1(\RR^d)$, where
$(\mu_{\Sigma,\Sigma'})$ is the affine Gaussian scale space with constant $b$
(\ref{def:affine-cascade-family}, \ref{eq:affine-gaussian}). The constant $b$ is then unique.
\statusT\quad\emph{(Cites nothing: it assembles \ref{lem:operators-to-kernels}(4),
\ref{thm:affine-rigidity} and \ref{lem:kernels-to-operators}(2). It is \ref{thm:affine-rigidity} read
on the operator axioms of the case $(\GL(d), \text{Loewner})$, which are the axioms of
\ref{lem:operators-to-kernels}(4) and of the draft's Definition~2.4. \textbf{Not claimed}:
the case $d = 1$, as for \ref{thm:affine-rigidity}.)}
\end{theorem}

\begin{proof}
\leanok
\emph{Necessity.} By \ref{lem:operators-to-kernels}(4) the measures $\mu_{\Sigma,\Sigma'}$
with $\Phi_{\Sigma,\Sigma'}f = \mu_{\Sigma,\Sigma'} * f$ form an affine cascade family in dimension
$d$, and by \ref{thm:affine-rigidity} it is the affine Gaussian scale space with some $b > 0$.

\emph{Sufficiency.} By \ref{thm:affine-rigidity} the affine Gaussian scale space with constant
$b > 0$ is an affine cascade family, that is a $(\GL(d), \text{Loewner})$-cascade family of kernels,
and $(\GL(d), \text{Loewner})$ is a wedge datum (\ref{lem:gw-scheme-claims}(2)). By
\ref{lem:kernels-to-operators}(2) the operators $f \mapsto \mu_{\Sigma,\Sigma'} * f$ form a
$(\GL(d), \text{Loewner})$-cascade family of operators.

\emph{Uniqueness of $b$.} If $\Phi_{0,I}$ is convolution by the Gaussian laws of constants $b$ and
$b'$, the uniqueness clause of \ref{lem:operators-to-kernels}(1) makes the two laws equal, and their
transforms \ref{eq:affine-gaussian} (\ref{lem:affine-gaussian-transform}) at $\Sigma = 0$, $\Sigma' = I$
and $\omega = e_1$ give $e^{-b} = e^{-b'}$.
\end{proof}
